% =====================================================================
\chapter{Soluciones de los ejercicios}
\label{ap:soluciones}
% =====================================================================

Cada solución se presenta en cinco partes: datos, planteamiento, desarrollo, resultado y comprobación. En los pasos intermedios se redondea a cuatro cifras significativas salvo que se indique otra cosa; los valores de $\Phi$ y los cuantiles se toman de las tablas del \cref{ap:tablas} o se dan con cuatro decimales. En los ejercicios de demostración, el apartado «Desarrollo» contiene la demostración redactada tal como debería entregarse.

% ---------------------------------------------------------------------
\section{Tema 1}

\begin{solucion}{eje:t1-vf}
\planteamiento contrastar cada afirmación con las definiciones del \cref{tema:1}.
\pasos
\begin{enumerate}[label=(\alph*)]
  \item \textbf{Verdadera} (si $n\ge 3$). El mayor dato sigue siendo el mayor después de multiplicarlo por $10$, de modo que el orden de los demás no cambia y la mediana, que depende solo de los valores centrales, tampoco. (Con $n=2$ la mediana es la media de los dos datos y sí cambia.)
  \item \textbf{Falsa.} $s$ es la raíz cuadrada no negativa de una suma de cuadrados dividida por $n-1>0$.
  \item \textbf{Verdadera.} Si $x_i=c$ para todo $i$, entonces $\xbar=c$, todas las desviaciones son nulas y $s=0$; además $Q_1=Q_3=c$.
  \item \textbf{Falsa.} Por el \cref{teo:t1-r}, los puntos están sobre una recta si y solo si $\abs r=1$.
  \item \textbf{Verdadera.} Por la \cref{prop:t1-lineal} con $a=1$, $b=5$: $s_y^{2}=s_x^{2}$ y $\RIC_y=\RIC_x$.
\end{enumerate}
\resultado \textbf{(a) V (para $n\ge3$); (b) F; (c) V; (d) F; (e) V.}
\end{solucion}

\begin{solucion}{eje:t1-basico}
\datos $4,7,2,9,5,6,3,25,6$; $n=9$.
\planteamiento ordenar, aplicar las \cref{def:t1-centro,def:t1-cuartiles,def:t1-caja} y la fórmula de cálculo (\cref{prop:t1-abreviada}).
\pasos
\begin{enumerate}[label=\arabic*.]
  \item Ordenados: $2,3,4,5,6,6,7,9,25$. Suma $67$; $\xbar=67/9\approx 7{,}444$.
  \item Mediana: el dato quinto, $6$. Moda: $6$ (aparece dos veces).
  \item Mitad inferior $2,3,4,5$: $Q_1=(3+4)/2=3{,}5$. Mitad superior $6,7,9,25$: $Q_3=(7+9)/2=8$. $\RIC=4{,}5$.
  \item Vallas: $3{,}5-1{,}5\cdot4{,}5=-3{,}25$ y $8+6{,}75=14{,}75$. El dato $25$ es atípico.
  \item $\sum x_i^{2}=16+49+4+81+25+36+9+625+36=881$; $n\xbar^{2}=67^{2}/9=4489/9\approx 498{,}78$; suma de cuadrados $881-498{,}78=382{,}22$; $s^{2}=382{,}22/8\approx 47{,}78$; $s\approx 6{,}912$.
  \item Sin el $25$ ($n=8$): suma $42$, $\xbar=5{,}25$; mediana $(5+6)/2=5{,}5$; $\sum x_i^{2}=256$, suma de cuadrados $256-8\cdot5{,}25^{2}=256-220{,}5=35{,}5$, $s^{2}=35{,}5/7\approx5{,}071$, $s\approx 2{,}252$.
\end{enumerate}
\resultado \textbf{$\xbar\approx7{,}44$; mediana $6$; moda $6$; $Q_1=3{,}5$; $Q_3=8$; $\RIC=4{,}5$; vallas $-3{,}25$ y $14{,}75$; $s^{2}\approx47{,}78$; $s\approx6{,}91$; el $25$ es atípico.} Sin él, $\xbar=5{,}25$, mediana $5{,}5$, $s\approx2{,}25$: la media baja unas dos unidades y $s$ se reduce a menos de la tercera parte, mientras que la mediana apenas cambia.
\comprobacion la suma de desviaciones con todos los datos es $67-9\cdot\frac{67}{9}=0$, como exige la \cref{prop:t1-extremal}.
\end{solucion}

\begin{solucion}{eje:t1-edades}
\datos $n=20$ edades; $\sum x_i=210{,}5$; $\sum x_i^{2}=2225{,}25$.
\planteamiento aplicar las definiciones; los datos ya están ordenados.
\pasos
\begin{enumerate}[label=\arabic*.]
  \item $\xbar=210{,}5/20=10{,}525$.
  \item Mediana: media de los datos décimo y undécimo, ambos $10{,}5$; mediana $10{,}5$.
  \item Mitad inferior (diez primeros): sus valores quinto y sexto son $10$ y $10$, luego $Q_1=10$. Mitad superior (diez últimos): sus valores quinto y sexto son $11$ y $11$, luego $Q_3=11$.
  \item Suma de cuadrados: $2225{,}25-20\cdot10{,}525^{2}=2225{,}25-2215{,}5125=9{,}7375$; $s^{2}=9{,}7375/19\approx0{,}5125$; $s\approx0{,}7159$.
  \item Intervalo $\xbar\pm s=(9{,}809;\allowbreak\ 11{,}241)$. Contiene los valores $10$ (cuatro veces), $10{,}5$ (cuatro) y $11$ (seis): $14$ datos.
\end{enumerate}
\resultado \textbf{$\xbar=10{,}525$; mediana $10{,}5$; $Q_1=10$; $Q_3=11$; $s\approx0{,}716$; el $70\,\%$ de los datos está a menos de una desviación típica de la media.}
\comprobacion la fuente da $s\approx0{,}72$. El $70\,\%$ es parecido al $68\,\%$ que se esperaría con datos de forma aproximadamente normal (\cref{tema:3}).
\end{solucion}

\begin{solucion}{eje:t1-cajas}
\datos mañana $(32;56;74{,}5;82{,}5;99)$; tarde $(25{,}5;78;81;89;98)$.
\planteamiento rango $=$ máximo $-$ mínimo; $\RIC=Q_3-Q_1$; vallas de la \cref{def:t1-caja}.
\pasos
\begin{enumerate}[label=\arabic*.]
  \item Mañana: rango $99-32=67$; $\RIC=82{,}5-56=26{,}5$.
  \item Tarde: rango $98-25{,}5=72{,}5$; $\RIC=89-78=11$.
  \item La mitad central más dispersa es la de la mañana ($26{,}5>11$), aunque el rango de la tarde sea mayor.
  \item Vallas de la mañana: $56-1{,}5\cdot26{,}5=16{,}25$ y $82{,}5+39{,}75=122{,}25$; el mínimo $32$ no es atípico.
  \item Vallas de la tarde: $78-16{,}5=61{,}5$ y $89+16{,}5=105{,}5$; el mínimo $25{,}5<61{,}5$ es atípico, y el bigote inferior acaba en el menor dato no atípico, que no se conoce con el resumen (está entre $61{,}5$ y $78$).
  \item Esbozo: la caja de la tarde es estrecha y está desplazada a la derecha (de $78$ a $89$, mediana $81$), con el $25{,}5$ marcado aparte; la de la mañana es ancha (de $56$ a $82{,}5$, mediana $74{,}5$) con bigotes hasta $32$ y $99$.
\end{enumerate}
\resultado \textbf{Rangos $67$ y $72{,}5$; $\RIC$ $26{,}5$ y $11$; la mañana tiene la mitad central más dispersa; solo el mínimo de la tarde es atípico.}
\comprobacion el rango de la tarde es mayor solo por un dato: el $\RIC$, robusto, revela que el grupo de la tarde es más homogéneo.
\end{solucion}

\begin{solucion}{eje:t1-agrupados}
\datos marcas $m_j=1,4,7,10,13,16$; frecuencias $f_j=1,6,10,7,0,2$; $n=26$.
\planteamiento $\xbar\approx\frac1n\sum_j f_jm_j$ y $s^{2}\approx\frac{1}{n-1}\bigl(\sum_j f_jm_j^{2}-n\xbar^{2}\bigr)$.
\pasos
\begin{enumerate}[label=\arabic*.]
  \item $\sum f_jm_j=1+24+70+70+0+32=197$; $\xbar=197/26\approx7{,}577$.
  \item $\sum f_jm_j^{2}=1+96+490+700+0+512=1799$.
  \item $n\xbar^{2}=197^{2}/26=38\,809/26\approx1492{,}65$; suma de cuadrados $1799-1492{,}65=306{,}35$.
  \item $s^{2}\approx306{,}35/25\approx12{,}25$; $s\approx3{,}50$.
\end{enumerate}
\resultado \textbf{$\xbar\approx7{,}58$ y $s\approx3{,}50$.}
\comprobacion la fuente obtiene los mismos valores. Son aproximaciones: se ha supuesto que todos los datos de cada clase valen la marca de clase.
\end{solucion}

\begin{solucion}{eje:t1-abreviada}
\datos $n=10$, $\sum x_i=50$, $\sum x_i^{2}=290$; nuevo dato $x_{11}=5$.
\planteamiento fórmula de cálculo (\cref{prop:t1-abreviada}); las sumas se actualizan sumando el nuevo dato.
\pasos
\begin{enumerate}[label=\arabic*.]
  \item $\xbar=50/10=5$; suma de cuadrados $290-10\cdot25=40$; $s^{2}=40/9\approx4{,}444$.
  \item Con el nuevo dato: $\sum x_i=55$, $\sum x_i^{2}=290+25=315$, $n=11$, $\xbar=5$.
  \item Suma de cuadrados $315-11\cdot25=40$; $s^{2}=40/10=4$.
\end{enumerate}
\resultado \textbf{(a) $\xbar=5$, $s^{2}=40/9\approx4{,}44$. (b) $\xbar=5$, $s^{2}=4$.}
\comprobacion el dato añadido coincide con la media, de modo que no altera la suma de cuadrados ($40$); solo aumenta el denominador, y la varianza disminuye.
\end{solucion}

\begin{solucion}{eje:t1-celsius}
\datos $18,21,25,19,22,30,24$ ($n=7$); $F=1{,}8C+32$.
\planteamiento calcular en Celsius y transformar con la \cref{prop:t1-lineal} ($a=1{,}8>0$).
\pasos
\begin{enumerate}[label=\arabic*.]
  \item Ordenados: $18,19,21,22,24,25,30$; suma $159$; $\xbar=159/7\approx22{,}71$; mediana $22$.
  \item Mitad inferior $18,19,21$: $Q_1=19$; superior $24,25,30$: $Q_3=25$; $\RIC=6$.
  \item $\sum x_i^{2}=3711$; suma de cuadrados $3711-159^{2}/7=3711-3611{,}57=99{,}43$; $s^{2}=99{,}43/6\approx16{,}57$; $s\approx4{,}071$.
  \item Fahrenheit: $\xbar_F=1{,}8\cdot22{,}714+32\approx72{,}89$; mediana $1{,}8\cdot22+32=71{,}6$; $s_F=1{,}8\cdot4{,}071\approx7{,}327$; $\RIC_F=1{,}8\cdot6=10{,}8$.
  \item Las puntuaciones $z$ no cambian: $z_i^{F}=\dfrac{1{,}8x_i+32-(1{,}8\xbar+32)}{1{,}8\,s}=\dfrac{x_i-\xbar}{s}$.
\end{enumerate}
\resultado \textbf{Celsius: $\xbar\approx22{,}71$, mediana $22$, $s\approx4{,}07$, $\RIC=6$. Fahrenheit: $\xbar\approx72{,}89$, mediana $71{,}6$, $s\approx7{,}33$, $\RIC=10{,}8$. Las puntuaciones $z$ son las mismas.}
\comprobacion el desplazamiento $+32$ afecta a la media y a la mediana, pero no a $s$ ni al $\RIC$.
\end{solucion}

\begin{solucion}{eje:t1-zscore}
\datos Juan: $x=2{,}85$, media $3{,}0$, $s=0{,}7$. Alí: $x=77$, media $80$, $s=10$.
\planteamiento comparar puntuaciones tipificadas (\cref{def:t1-z}).
\pasos
\begin{enumerate}[label=\arabic*.]
  \item $z_{\text{Juan}}=(2{,}85-3{,}0)/0{,}7=-0{,}15/0{,}7\approx-0{,}214$.
  \item $z_{\text{Alí}}=(77-80)/10=-0{,}3$.
\end{enumerate}
\resultado \textbf{Juan}, porque su nota está a $0{,}21$ desviaciones típicas por debajo de la media de su centro, frente a $0{,}30$ en el caso de Alí.
\comprobacion la comparación directa de $2{,}85$ con $77$ carece de sentido porque las escalas son distintas; la tipificación las hace comparables.
\end{solucion}

\begin{solucion}{eje:t1-cov-tabla}
\datos $X,Y\in\{-1,1\}$ con $\PP(X=1)=\PP(Y=1)=1/2$; $c=\PP(X=1,Y=1)$.
\planteamiento completar la tabla con las marginales y usar $\Cov(X,Y)=\E[XY]-\E[X]\E[Y]$.
\pasos
\begin{enumerate}[label=\arabic*.]
  \item Como $\PP(X=1)=1/2$, $\PP(X=1,Y=-1)=1/2-c$; análogamente $\PP(X=-1,Y=1)=1/2-c$ y $\PP(X=-1,Y=-1)=1-c-2(1/2-c)=c$. Debe ser $0\le c\le1/2$.
  \item $\E[X]=\E[Y]=0$ y $\Var X=\Var Y=\E[X^{2}]=1$.
  \item $XY=1$ con probabilidad $2c$ y $XY=-1$ con probabilidad $1-2c$, luego $\E[XY]=2c-(1-2c)=4c-1$.
  \item $\Cov(X,Y)=4c-1$ y $\Cor(X,Y)=(4c-1)/\sqrt{1\cdot1}=4c-1$.
  \item Independencia: debe cumplirse $\PP(X=1,Y=1)=\frac12\cdot\frac12$, es decir, $c=1/4$; entonces las cuatro casillas valen $1/4$ y la tabla es producto de las marginales.
  \item $\Cor=1\iff c=1/2$ (entonces $Y=X$); $\Cor=-1\iff c=0$ (entonces $Y=-X$).
\end{enumerate}
\resultado \textbf{$\Cov(X,Y)=\Cor(X,Y)=4c-1$; independientes si y solo si $c=1/4$; correlación $1$ si $c=1/2$ y $-1$ si $c=0$.}
\comprobacion $c\in[0,1/2]$ da $\Cor\in[-1,1]$, de acuerdo con la acotación de la correlación.
\end{solucion}

\begin{solucion}{eje:t1-agrupar}
\datos $n_1=10$, $\xbar_1=70$, $s_1=8$; $n_2=15$, $\xbar_2=80$, $s_2=6$.
\planteamiento la suma de cuadrados respecto de la media global $\xbar$ se descompone por la \cref{prop:t1-extremal}(ii): $\sum_{\text{grupo }g}(x_i-\xbar)^{2}=(n_g-1)s_g^{2}+n_g(\xbar_g-\xbar)^{2}$.
\pasos
\begin{enumerate}[label=\arabic*.]
  \item $\xbar=(10\cdot70+15\cdot80)/25=(700+1200)/25=76$.
  \item Grupo 1: $9\cdot64+10\cdot(70-76)^{2}=576+360=936$.
  \item Grupo 2: $14\cdot36+15\cdot(80-76)^{2}=504+240=744$.
  \item Total: $936+744=1680$; $s^{2}=1680/24=70$; $s=\sqrt{70}\approx8{,}367$.
\end{enumerate}
\resultado \textbf{$\xbar=76$ y $s\approx8{,}37$.}
\comprobacion $s$ es mayor que las dos desviaciones de grupo porque se añade la variabilidad entre las medias de los grupos.
\end{solucion}

\begin{solucion}{eje:t1-mediana-min}
\planteamiento estudiar la pendiente de $g$ en cada intervalo entre datos ordenados consecutivos.
\pasos
\begin{enumerate}[label=\arabic*.]
  \item $g$ es continua y lineal a trozos. Sea $a$ un punto que no coincide con ningún dato y sean $L(a)=\#\{i:x_i<a\}$ y $U(a)=\#\{i:x_i>a\}$, de modo que $L(a)+U(a)=n$. Cerca de $a$, cada término con $x_i<a$ vale $a-x_i$ (pendiente $+1$) y cada término con $x_i>a$ vale $x_i-a$ (pendiente $-1$). Por tanto, $g'(a)=L(a)-U(a)$.
  \item Si $L(a)<n/2$, entonces $U(a)>n/2$ y $g'(a)<0$: $g$ decrece. Si $L(a)>n/2$, $g$ crece.
  \item Caso $n=2m+1$ impar. Para $a<x_{(m+1)}$ (y fuera de los datos) hay a lo sumo $m<n/2$ datos menores que $a$, y $g$ decrece; para $a>x_{(m+1)}$ hay al menos $m+1>n/2$, y $g$ crece. Por continuidad, el mínimo está en $a=x_{(m+1)}$, la mediana.
  \item Caso $n=2m$ par. Para $a<x_{(m)}$, $L(a)\le m-1<m$ y $g$ decrece; para $a>x_{(m+1)}$, $L(a)\ge m+1$ y $g$ crece; para $x_{(m)}<a<x_{(m+1)}$, $L(a)=U(a)=m$ y $g$ es constante. Luego $g$ es mínima en todo el intervalo $[x_{(m)},x_{(m+1)}]$, que contiene la mediana $\tfrac12(x_{(m)}+x_{(m+1)})$.
\end{enumerate}
\resultado \textbf{$g$ alcanza su mínimo en la mediana} (si $n$ es par, en todo el intervalo entre los dos datos centrales). $\blacksquare$
\comprobacion con los datos $1,2,10$: $g(2)=1+0+8=9$, $g(3)=2+1+7=10$ y $g(\xbar)=g(13/3)\approx3{,}33+2{,}33+5{,}67=11{,}33$; la media no minimiza la suma de valores absolutos.
\end{solucion}

\begin{solucion}{eje:t1-r-invariante}
\planteamiento por la \cref{prop:t1-lineal}, $x_i'-\xbar'=a(x_i-\xbar)$ e $y_i'-\ybar'=c(y_i-\ybar)$.
\pasos
\begin{enumerate}[label=\arabic*.]
  \item $s_{x'y'}=\frac{1}{n-1}\sum_i a(x_i-\xbar)\,c(y_i-\ybar)=ac\,s_{xy}$.
  \item $s_{x'}=\abs a\,s_x$ y $s_{y'}=\abs c\,s_y$.
  \item $r_{x'y'}=\dfrac{ac\,s_{xy}}{\abs a\abs c\,s_xs_y}=\dfrac{ac}{\abs{ac}}\,r_{xy}=\operatorname{signo}(ac)\,r_{xy}$.
  \item Un cambio de unidades es una transformación con $a>0$ (o $c>0$), luego no altera $r$.
\end{enumerate}
\resultado \textbf{$s_{x'y'}=ac\,s_{xy}$ y $r_{x'y'}=\operatorname{signo}(ac)\,r_{xy}$}; la correlación es adimensional. $\blacksquare$
\comprobacion con los datos del \cref{ej:t1-r} y $x'=10x$: la covarianza pasa de $1{,}5$ a $15$ y $r$ sigue valiendo $0{,}775$.
\end{solucion}

\begin{solucion}{eje:t1-chebyshev}
\planteamiento acotar inferiormente la suma de cuadrados por la contribución de los datos alejados.
\pasos
\begin{enumerate}[label=\arabic*.]
  \item Sea $A=\{i:\abs{x_i-\xbar}\ge ks\}$ y $N_A=\abs A$.
  \item Como todos los sumandos son no negativos,
  \[
  (n-1)s^{2}=\sum_{i=1}^{n}(x_i-\xbar)^{2}\ge\sum_{i\in A}(x_i-\xbar)^{2}\ge N_A\,k^{2}s^{2}.
  \]
  \item Dividiendo por $nk^{2}s^{2}>0$: $N_A/n\le(n-1)/(nk^{2})<1/k^{2}$.
  \item Con $k=2$: menos del $25\,\%$ de los datos dista $2s$ o más de la media; al menos el $75\,\%$ está en $(\xbar-2s,\xbar+2s)$.
\end{enumerate}
\resultado \textbf{La fracción es a lo sumo $(n-1)/(nk^{2})<1/k^{2}$}; para $k=2$, al menos el $75\,\%$ de los datos está a menos de dos desviaciones típicas de la media, sea cual sea la forma de los datos. $\blacksquare$
\comprobacion en el \cref{ej:t1-tokens}, $\xbar\pm2s=(-1{,}45;\allowbreak\ 36{,}2)$ deja fuera solo el $40$: $1/8=12{,}5\,\%\le7/32\approx21{,}9\,\%$.
\end{solucion}

\begin{solucion}{eje:t1-welford}
\planteamiento escribir $m\xbar_m=(m-1)\xbar_{m-1}+x_m$ y aplicar la \cref{prop:t1-extremal}(ii) a los $m-1$ primeros datos con $a=\xbar_m$.
\pasos
\begin{enumerate}[label=\arabic*.]
  \item Media: $\xbar_m=\dfrac{(m-1)\xbar_{m-1}+x_m}{m}=\xbar_{m-1}+\dfrac{x_m-\xbar_{m-1}}{m}$.
  \item Por la \cref{prop:t1-extremal}(ii), $\sum_{i=1}^{m-1}(x_i-\xbar_m)^{2}=M_{m-1}+(m-1)(\xbar_{m-1}-\xbar_m)^{2}$. Sumando el término $i=m$: $M_m=M_{m-1}+(m-1)(\xbar_{m-1}-\xbar_m)^{2}+(x_m-\xbar_m)^{2}$.
  \item Sea $d=x_m-\xbar_{m-1}$. Por el paso 1, $\xbar_m-\xbar_{m-1}=d/m$ y $x_m-\xbar_m=d-d/m=d(m-1)/m$.
  \item Sustituyendo: $(m-1)\dfrac{d^{2}}{m^{2}}+\dfrac{(m-1)^{2}d^{2}}{m^{2}}=\dfrac{(m-1)d^{2}(1+m-1)}{m^{2}}=\dfrac{(m-1)d^{2}}{m}=d\cdot\dfrac{(m-1)d}{m}=(x_m-\xbar_{m-1})(x_m-\xbar_m)$.
  \item Datos $2,4,9$: $\xbar_1=2$, $M_1=0$; $\xbar_2=2+2/2=3$, $M_2=0+(4-2)(4-3)=2$; $\xbar_3=3+6/3=5$, $M_3=2+(9-3)(9-5)=26$.
\end{enumerate}
\resultado \textbf{Las dos fórmulas son ciertas}; con $2,4,9$ se obtiene $\xbar_3=5$, $M_3=26$ y $s^{2}=26/2=13$. $\blacksquare$
\comprobacion directamente: $(2-5)^{2}+(4-5)^{2}+(9-5)^{2}=9+1+16=26$.
\end{solucion}

\begin{solucion}{eje:t1-varianza-pares}
\planteamiento desarrollar la suma sobre todos los pares ordenados y usar $\sum_i(x_i-\xbar)=0$.
\pasos
\begin{enumerate}[label=\arabic*.]
  \item Sean $u_i=x_i-\xbar$; como $x_i-x_j=u_i-u_j$, se trabaja con los $u_i$, que suman cero.
  \item Sumando sobre todos los pares ordenados $(i,j)$, incluidos $i=j$ (que aportan $0$):
  \[
  \sum_{i=1}^{n}\sum_{j=1}^{n}(u_i-u_j)^{2}=\sum_{i,j}u_i^{2}+\sum_{i,j}u_j^{2}-2\Bigl(\sum_iu_i\Bigr)\Bigl(\sum_ju_j\Bigr)=2n\sum_iu_i^{2}-0.
  \]
  \item Cada par no ordenado $\{i,j\}$ con $i\neq j$ aparece dos veces, luego $\sum_{i<j}(x_i-x_j)^{2}=n\sum_iu_i^{2}=n(n-1)s^{2}$.
  \item Dividiendo por $n(n-1)$ se obtiene la identidad.
\end{enumerate}
\resultado \textbf{$s^{2}=\frac{1}{n(n-1)}\sum_{i<j}(x_i-x_j)^{2}$.} $\blacksquare$
\comprobacion con $2,4,9$: $(2-4)^{2}+(2-9)^{2}+(4-9)^{2}=4+49+25=78$ y $78/(3\cdot2)=13$, que coincide con el \cref{eje:t1-welford}.
\end{solucion}

% ---------------------------------------------------------------------
\section{Tema 2}

\begin{solucion}{eje:t2-conceptos}
\planteamiento aplicar las \cref{def:t2-mas,def:t2-parametro,def:t2-estimador}.
\pasos
\begin{enumerate}[label=\arabic*.]
  \item Población: el banco de un millón de preguntas (o, mejor, la variable «acierto» en cada una de ellas).
  \item Muestra: las $500$ preguntas elegidas al azar, con variables $X_1,\dots,X_{500}$, donde $X_i=1$ si el modelo acierta la pregunta $i$.
  \item Parámetro: la exactitud real $p$, fracción de aciertos en todo el banco; es un número fijo y desconocido.
  \item Estimador: $\widehat p=\Xbar=\frac{1}{500}\sum_iX_i$, una variable aleatoria. Estimación: su valor observado $410/500=0{,}82$.
\end{enumerate}
\resultado \textbf{El número $0{,}82$ no es aleatorio (es un valor observado); $\widehat p$ sí lo es}, porque con otra muestra de $500$ preguntas tomaría otro valor.
\comprobacion la muestra ($500$) es mucho menor que la población ($10^{6}$), de modo que el muestreo sin reemplazamiento se puede tratar como \iid\ (\cref{def:t2-mas}).
\end{solucion}

\begin{solucion}{eje:t2-enumeracion}
\datos población $\{1,3,5\}$ equiprobable; $n=2$ con reemplazamiento.
\planteamiento enumerar las $3^{2}=9$ muestras equiprobables y calcular $\xbar$ y $s^{2}=(x_1-x_2)^{2}/2$ en cada una.
\pasos
\begin{enumerate}[label=\arabic*.]
  \item $\mu=(1+3+5)/3=3$; $\sigma^{2}=\bigl((1-3)^{2}+0+(5-3)^{2}\bigr)/3=8/3$.
  \item Muestras y medias: $(1,1)\to1$; $(1,3),(3,1)\to2$; $(1,5),(3,3),(5,1)\to3$; $(3,5),(5,3)\to4$; $(5,5)\to5$. Distribución de $\Xbar$: valores $1,2,3,4,5$ con probabilidades $\frac19,\frac29,\frac39,\frac29,\frac19$.
  \item $\E[\Xbar]=(1+4+9+8+5)/9=27/9=3=\mu$.
  \item $\Var(\Xbar)=\bigl(4\cdot1+1\cdot2+0\cdot3+1\cdot2+4\cdot1\bigr)/9=12/9=4/3=\sigma^{2}/2$.
  \item Para $n=2$, $s^{2}=\bigl((x_1-\xbar)^{2}+(x_2-\xbar)^{2}\bigr)/1=(x_1-x_2)^{2}/2$. Vale $0$ en las tres muestras con $x_1=x_2$, $2$ en las cuatro con $\abs{x_1-x_2}=2$ y $8$ en las dos con $\abs{x_1-x_2}=4$.
  \item $\E[S^{2}]=(0\cdot3+2\cdot4+8\cdot2)/9=24/9=8/3=\sigma^{2}$.
  \item $\widehat\sigma^{2}=S^{2}/2$, luego $\E[\widehat\sigma^{2}]=4/3$.
\end{enumerate}
\resultado \textbf{$\mu=3$, $\sigma^{2}=8/3$; $\E[\Xbar]=3$, $\Var(\Xbar)=4/3$; $\E[S^{2}]=8/3$; $\E[\widehat\sigma^{2}]=4/3$}, la mitad de $\sigma^{2}$.
\comprobacion el \cref{cor:t2-sesgo} predice $\E[\widehat\sigma^{2}]=\frac{n-1}{n}\sigma^{2}=\frac12\cdot\frac83=\frac43$.
\end{solucion}

\begin{solucion}{eje:t2-tres-estimadores}
\datos $X_1,X_2,X_3$ \iid, media $\mu$, varianza $\sigma^{2}$.
\planteamiento para $T=\sum_ia_iX_i$: $\E[T]=\mu\sum_ia_i$ y $\Var(T)=\sigma^{2}\sum_ia_i^{2}$ (\cref{prop:t2-props}); después, \cref{prop:t2-ecm}.
\pasos
\begin{enumerate}[label=\arabic*.]
  \item $T_1$: $a=(\frac13,\frac13,\frac13)$; $\sum a_i=1$, sesgo $0$; $\sum a_i^{2}=\frac13$; $\ECM=\Var=\sigma^{2}/3$.
  \item $T_2$: $a=(\frac14,\frac12,\frac14)$; $\sum a_i=1$, sesgo $0$; $\sum a_i^{2}=\frac1{16}+\frac14+\frac1{16}=\frac38$; $\ECM=3\sigma^{2}/8$.
  \item $T_3$: $a=(1,1,-1)$; $\sum a_i=1$, sesgo $0$; $\sum a_i^{2}=3$; $\ECM=3\sigma^{2}$.
  \item $T_4$: $a=(\frac14,\frac14,\frac14)$; $\E[T_4]=\frac34\mu$, sesgo $-\mu/4$; $\Var=\frac3{16}\sigma^{2}$; $\ECM=\frac{3\sigma^{2}}{16}+\frac{\mu^{2}}{16}$.
  \item Entre los insesgados, $\frac13<\frac38<3$: el mejor es $T_1$.
  \item $\ECM(T_4)<\ECM(T_1)\iff\frac{3\sigma^{2}+\mu^{2}}{16}<\frac{\sigma^{2}}{3}\iff\mu^{2}<\frac{16\sigma^{2}}{3}-3\sigma^{2}=\frac{7\sigma^{2}}{3}$.
\end{enumerate}
\resultado \textbf{Sesgos $0,0,0,-\mu/4$; varianzas $\sigma^{2}/3$, $3\sigma^{2}/8$, $3\sigma^{2}$, $3\sigma^{2}/16$; el mejor insesgado es $\Xbar$; $T_4$ es mejor que $\Xbar$ si y solo si $\abs\mu<\sigma\sqrt{7/3}\approx1{,}528\,\sigma$.}
\comprobacion para $\mu=0$, $T_4$ es insesgado y tiene varianza $3\sigma^{2}/16<\sigma^{2}/3$: encoger hacia el cero ayuda cuando el parámetro está cerca de cero.
\end{solucion}

\begin{solucion}{eje:t2-poisson}
\datos $X_i\sim\Pois(\lambda)$, $\E X_i=\Var X_i=\lambda$; $n=50$, $\xbar=3{,}2$.
\planteamiento \cref{teo:t2-media} con $\mu=\sigma^{2}=\lambda$.
\pasos
\begin{enumerate}[label=\arabic*.]
  \item $\E[\widehat\lambda]=\lambda$: insesgado.
  \item $\Var(\widehat\lambda)=\lambda/n$; $\ee(\widehat\lambda)=\sqrt{\lambda/n}$; $\ECM=\Var=\lambda/n$.
  \item Estimación: $\widehat\lambda=3{,}2$. Error estándar estimado: $\sqrt{3{,}2/50}=\sqrt{0{,}064}\approx0{,}2530$.
\end{enumerate}
\resultado \textbf{$\widehat\lambda$ es insesgado, con $\Var=\ECM=\lambda/n$ y $\ee=\sqrt{\lambda/n}$; con los datos, $\widehat\lambda=3{,}2$ y $\widehat{\ee}\approx0{,}253$.}
\comprobacion el error estándar es pequeño frente a $3{,}2$ (un $8\,\%$), coherente con $n=50$ bloques.
\end{solucion}

\begin{solucion}{eje:t2-expo-sesgo}
\datos $\E[1/\Xbar]=n\lambda/(n-1)$; $n=5$; $\sum x_i=10$.
\planteamiento sesgo $=\E[\widehat\lambda]-\lambda$; elegir $c$ para que $\E[c/\Xbar]=\lambda$.
\pasos
\begin{enumerate}[label=\arabic*.]
  \item $\sesgo(1/\Xbar)=\dfrac{n\lambda}{n-1}-\lambda=\dfrac{\lambda}{n-1}>0$: sobrestima $\lambda$.
  \item $\E[c/\Xbar]=c\,n\lambda/(n-1)=\lambda\iff c=(n-1)/n$. El estimador insesgado es $\dfrac{n-1}{n\Xbar}=\dfrac{n-1}{\sum_iX_i}$.
  \item Datos: $\xbar=2$; $1/\xbar=0{,}5$; insesgado: $4/10=0{,}4$.
\end{enumerate}
\resultado \textbf{Sesgo $\lambda/(n-1)$; estimador insesgado $(n-1)/\sum_iX_i$; estimaciones $0{,}5$ y $0{,}4$.}
\comprobacion el sesgo tiende a cero cuando $n\to\infty$ y la corrección $(n-1)/n\to1$; con $n=5$ la diferencia es del $20\,\%$.
\end{solucion}

\begin{solucion}{eje:t2-sumas-solapadas}
\datos $\Var(X_i)=1/4$; $X=\sum_{i=1}^{n}X_i$; $Y=\sum_{i=n-7}^{2n-8}X_i$.
\planteamiento bilinealidad de la covarianza (\cref{prop:t2-props}(v)): $\Cov(X_i,X_j)=0$ si $i\ne j$ y $\Var(X_i)$ si $i=j$; solo cuentan los índices comunes.
\pasos
\begin{enumerate}[label=\arabic*.]
  \item $Y$ tiene $(2n-8)-(n-7)+1=n$ sumandos, luego $\Var(Y)=n/4=\Var(X)$.
  \item Índices comunes: de $n-7$ a $n$, que son $8$ (como $n\ge8$, $n-7\ge1$).
  \item $\Cov(X,Y)=\sum_{i=n-7}^{n}\Var(X_i)=8\cdot\frac14=2$.
  \item $\Cor(X,Y)=\dfrac{2}{\sqrt{(n/4)(n/4)}}=\dfrac{2}{n/4}=\dfrac8n$.
\end{enumerate}
\resultado \textbf{$\Cov(X,Y)=2$ y $\Cor(X,Y)=8/n$.}
\comprobacion para $n=8$ se obtiene $\Cor=1$: entonces $Y=X_1+\dots+X_8=X$.
\end{solucion}

\begin{solucion}{eje:t2-benchmark}
\datos $n=400$; $300$ aciertos.
\planteamiento \cref{ej:t2-proporcion}: $\widehat{\ee}=\sqrt{\widehat p(1-\widehat p)/n}$.
\pasos
\begin{enumerate}[label=\arabic*.]
  \item $\widehat p=300/400=0{,}75$.
  \item $\widehat{\ee}=\sqrt{0{,}75\cdot0{,}25/400}=\sqrt{0{,}000\,468\,75}\approx0{,}02165$.
  \item El error estándar es proporcional a $1/\sqrt n$; para dividirlo por $2$ se necesita $4n=1600$ preguntas.
\end{enumerate}
\resultado \textbf{$\widehat p=0{,}75$, $\widehat{\ee}\approx0{,}0217$ (unos $2{,}2$ puntos porcentuales); $1600$ preguntas.}
\comprobacion $\sqrt{0{,}1875/1600}\approx0{,}01083$, la mitad de $0{,}02165$.
\end{solucion}

\begin{solucion}{eje:t2-combinar}
\datos $\E A=\E B=\theta$; $A$, $B$ independientes; $\Var A=4$, $\Var B=1$.
\planteamiento linealidad para el sesgo; varianza de una combinación de independientes; minimizar un polinomio de grado~2.
\pasos
\begin{enumerate}[label=\arabic*.]
  \item $\E[T_w]=w\theta+(1-w)\theta=\theta$ para todo $w$.
  \item $V(w)=\Var(T_w)=4w^{2}+(1-w)^{2}$. $V'(w)=8w-2(1-w)=10w-2=0\iff w=0{,}2$; $V''=10>0$, mínimo.
  \item $V(0{,}2)=4\cdot0{,}04+0{,}64=0{,}8$.
  \item En general, $V(w)=w^{2}\sigma_A^{2}+(1-w)^{2}\sigma_B^{2}$ y $V'(w)=0$ da $w^{*}=\dfrac{\sigma_B^{2}}{\sigma_A^{2}+\sigma_B^{2}}$, con $V(w^{*})=\dfrac{\sigma_A^{2}\sigma_B^{2}}{\sigma_A^{2}+\sigma_B^{2}}$.
\end{enumerate}
\resultado \textbf{$w^{*}=0{,}2$ y $\Var_{\min}=0{,}8$}; en general, cada estimador se pondera de forma inversamente proporcional a su varianza y $1/\Var_{\min}=1/\sigma_A^{2}+1/\sigma_B^{2}$.
\comprobacion $1/4+1/1=1{,}25=1/0{,}8$. La combinación es mejor que el mejor de los dos ($0{,}8<1$).
\end{solucion}

\begin{solucion}{eje:t2-chebyshev-n}
\datos $\varepsilon=0{,}05$; probabilidad máxima $0{,}05$; $\Var(\widehat p)=p(1-p)/n\le1/(4n)$.
\planteamiento \cref{lem:t2-chebyshev} con $W=\widehat p$.
\pasos
\begin{enumerate}[label=\arabic*.]
  \item $\PP(\abs{\widehat p-p}\ge0{,}05)\le\dfrac{p(1-p)}{n\cdot0{,}0025}\le\dfrac{1}{4n\cdot0{,}0025}=\dfrac{100}{n}$.
  \item $100/n\le0{,}05\iff n\ge2000$.
\end{enumerate}
\resultado \textbf{$n=2000$ preguntas bastan}, sea cual sea $p$.
\comprobacion la aproximación normal del \cref{tema:3} dará $n=385$ (\cref{eje:t3-tamano}); Chebyshev es válida para cualquier distribución, pero es muy conservadora.
\end{solucion}

\begin{solucion}{eje:t2-ecm-normal}
\datos $\E[S^{2}]=\sigma^{2}$; $\Var(S^{2})=2\sigma^{4}/(n-1)$; $\widehat\sigma^{2}=\frac{n-1}{n}S^{2}$.
\planteamiento \cref{prop:t2-ecm} para cada estimador y comparación de las dos expresiones.
\pasos
\begin{enumerate}[label=\arabic*.]
  \item $S^{2}$ es insesgado, luego $\ECM(S^{2})=\dfrac{2\sigma^{4}}{n-1}$.
  \item $\sesgo(\widehat\sigma^{2})=-\sigma^{2}/n$ (\cref{cor:t2-sesgo}) y $\Var(\widehat\sigma^{2})=\dfrac{(n-1)^{2}}{n^{2}}\cdot\dfrac{2\sigma^{4}}{n-1}=\dfrac{2(n-1)\sigma^{4}}{n^{2}}$.
  \item $\ECM(\widehat\sigma^{2})=\dfrac{2(n-1)\sigma^{4}}{n^{2}}+\dfrac{\sigma^{4}}{n^{2}}=\dfrac{(2n-1)\sigma^{4}}{n^{2}}$.
  \item Comparación: $\dfrac{2n-1}{n^{2}}<\dfrac{2}{n-1}\iff(2n-1)(n-1)<2n^{2}\iff2n^{2}-3n+1<2n^{2}\iff1<3n$, que se cumple para todo $n\ge1$.
\end{enumerate}
\resultado \textbf{$\ECM(S^{2})=2\sigma^{4}/(n-1)$ y $\ECM(\widehat\sigma^{2})=(2n-1)\sigma^{4}/n^{2}$, que es menor para todo $n\ge2$.} $\blacksquare$
\comprobacion con $n=2$: $\ECM(S^{2})=2\sigma^{4}$ y $\ECM(\widehat\sigma^{2})=\frac34\sigma^{4}$. La insesgadez no garantiza el menor error cuadrático medio.
\end{solucion}

\begin{solucion}{eje:t2-uniforme}
\datos $X_i\sim\Unif(0,\theta)$: densidad $1/\theta$ en $[0,\theta]$, media $\theta/2$, varianza $\theta^{2}/12$.
\planteamiento (a) \cref{teo:t2-media}; (b) la función de distribución del máximo por independencia; (c) comparar.
\pasos
\begin{enumerate}[label=\arabic*.]
  \item $\E[2\Xbar]=2\cdot\theta/2=\theta$; $\Var(2\Xbar)=4\cdot\dfrac{\theta^{2}}{12n}=\dfrac{\theta^{2}}{3n}$.
  \item $M\le x$ si y solo si todos los $X_i\le x$; por independencia, $\PP(M\le x)=\prod_i\PP(X_i\le x)=(x/\theta)^{n}$ para $0\le x\le\theta$. Derivando, la densidad de $M$ es $f_M(x)=nx^{n-1}/\theta^{n}$ en $[0,\theta]$.
  \item $\E[M]=\displaystyle\int_0^{\theta}x\,\frac{nx^{n-1}}{\theta^{n}}\,dx=\frac{n}{\theta^{n}}\cdot\frac{\theta^{n+1}}{n+1}=\frac{n\theta}{n+1}$; análogamente $\E[M^{2}]=\dfrac{n}{\theta^{n}}\cdot\dfrac{\theta^{n+2}}{n+2}=\dfrac{n\theta^{2}}{n+2}$.
  \item $\E[T_2]=\frac{n+1}{n}\cdot\frac{n\theta}{n+1}=\theta$: insesgado.
  \item $\Var(M)=\dfrac{n\theta^{2}}{n+2}-\dfrac{n^{2}\theta^{2}}{(n+1)^{2}}=\dfrac{n\theta^{2}\bigl[(n+1)^{2}-n(n+2)\bigr]}{(n+2)(n+1)^{2}}=\dfrac{n\theta^{2}}{(n+2)(n+1)^{2}}$.
  \item $\Var(T_2)=\dfrac{(n+1)^{2}}{n^{2}}\cdot\dfrac{n\theta^{2}}{(n+2)(n+1)^{2}}=\dfrac{\theta^{2}}{n(n+2)}$.
  \item $\dfrac{\theta^{2}}{n(n+2)}\le\dfrac{\theta^{2}}{3n}\iff3\le n+2\iff n\ge1$, con igualdad solo si $n=1$.
\end{enumerate}
\resultado \textbf{Ambos son insesgados; $\Var(T_1)=\theta^{2}/(3n)$ y $\Var(T_2)=\theta^{2}/(n(n+2))$.} $T_2$ es mejor para $n\ge2$ y su varianza decrece como $1/n^{2}$, no como $1/n$. $\blacksquare$
\comprobacion con $n=1$, $M=X_1$ y $T_2=2X_1=T_1$: deben coincidir, y en efecto $\theta^{2}/3=\theta^{2}/(1\cdot3)$.
\end{solucion}

\begin{solucion}{eje:t2-max-consistente}
\datos $\PP(M\le x)=(x/\theta)^{n}$; $\ECM(M)$ a partir del \cref{eje:t2-uniforme}.
\planteamiento (a) como $M\le\theta$, el suceso $\abs{M-\theta}>\varepsilon$ es $M<\theta-\varepsilon$; (b) calcular el $\ECM$.
\pasos
\begin{enumerate}[label=\arabic*.]
  \item Si $\varepsilon\ge\theta$, la probabilidad es $0$. Si $0<\varepsilon<\theta$: $\PP(\abs{M-\theta}>\varepsilon)=\PP(M<\theta-\varepsilon)=\bigl(1-\varepsilon/\theta\bigr)^{n}\to0$, porque $0<1-\varepsilon/\theta<1$.
  \item Sesgo: $\E[M]-\theta=-\theta/(n+1)\to0$. Varianza: $\dfrac{n\theta^{2}}{(n+2)(n+1)^{2}}\le\dfrac{\theta^{2}}{(n+1)^{2}}\to0$. Luego $\ECM(M)\to0$ y la \cref{prop:t2-ecm-consistencia} da la consistencia.
\end{enumerate}
\resultado \textbf{$M\convp\theta$}: el máximo es un estimador consistente (aunque sesgado). $\blacksquare$
\comprobacion con $\theta=1$, $\varepsilon=0{,}05$ y $n=100$: $0{,}95^{100}\approx0{,}0059$.
\end{solucion}

\begin{solucion}{eje:t2-cov-desv}
\planteamiento bilinealidad de la covarianza y $\Cov(X_j,X_i)=\sigma^{2}$ si $j=i$ y $0$ si $j\neq i$.
\pasos
\begin{enumerate}[label=\arabic*.]
  \item $\Cov(\Xbar,X_i)=\frac1n\sum_{j=1}^{n}\Cov(X_j,X_i)=\frac1n\sigma^{2}$.
  \item $\Cov(\Xbar,\Xbar)=\Var(\Xbar)=\sigma^{2}/n$.
  \item $\Cov(\Xbar,X_i-\Xbar)=\Cov(\Xbar,X_i)-\Cov(\Xbar,\Xbar)=\sigma^{2}/n-\sigma^{2}/n=0$.
\end{enumerate}
\resultado \textbf{$\Cov(\Xbar,X_i-\Xbar)=0$ para todo $i$.} $\blacksquare$
\comprobacion la incorrelación no implica en general independencia; para datos normales, el vector $(\Xbar,X_1-\Xbar,\dots,X_n-\Xbar)$ es normal multivariante y en ese caso sí la implica, lo que lleva a la independencia de $\Xbar$ y $S^{2}$ (\cref{teo:t3-fisher}).
\end{solucion}

\begin{solucion}{eje:t2-s-sesgo}
\planteamiento $\Var(S)=\E[S^{2}]-(\E S)^{2}\ge0$ y $\E[S^{2}]=\sigma^{2}$ (\cref{teo:t2-s2}).
\pasos
\begin{enumerate}[label=\arabic*.]
  \item $0\le\Var(S)=\sigma^{2}-(\E S)^{2}$, luego $(\E S)^{2}\le\sigma^{2}$.
  \item Como $S\ge0$, $\E S\ge0$, y tomando raíces: $\E S\le\sigma$.
  \item Igualdad $\iff\Var(S)=0\iff S$ es constante con probabilidad~1.
\end{enumerate}
\resultado \textbf{$\E[S]\le\sigma$}, con igualdad solo en el caso degenerado; $S$ subestima $\sigma$ en promedio. $\blacksquare$
\comprobacion con la población del \cref{eje:t2-enumeracion}: $S=\sqrt{S^{2}}$ vale $0$, $\sqrt2$ y $\sqrt8$ con probabilidades $3/9$, $4/9$ y $2/9$, luego $\E S=(4\sqrt2+2\sqrt8)/9=8\sqrt2/9\approx1{,}257<\sigma=\sqrt{8/3}\approx1{,}633$.
\end{solucion}

\begin{solucion}{eje:t2-passk}
\datos $n$ soluciones independientes, cada una correcta con probabilidad $p$; $C$ correctas.
\planteamiento introducir un subconjunto aleatorio $K$ de $k$ índices, elegido uniformemente entre los $\binom nk$ posibles e independiente de las soluciones, y aplicar la ley de la probabilidad total de dos formas.
\pasos
\begin{enumerate}[label=\arabic*.]
  \item Sea $A$ el suceso «todas las soluciones de índice en $K$ son incorrectas».
  \item Condicionando a los resultados de las $n$ soluciones (y, por tanto, a $C$), $K$ es un subconjunto uniforme, y el número de subconjuntos de $k$ índices formados solo por incorrectas es $\binom{n-C}{k}$ (es $0$ si $n-C<k$). Luego $\PP(A\mid\text{resultados})=\binom{n-C}{k}/\binom nk$.
  \item Tomando esperanzas (probabilidad total): $\PP(A)=\E\Bigl[\binom{n-C}{k}\big/\binom nk\Bigr]$.
  \item Condicionando ahora a $K=J$ para un subconjunto fijo $J$: las $k$ soluciones de $J$ son independientes y cada una es incorrecta con probabilidad $1-p$, luego $\PP(A\mid K=J)=(1-p)^{k}$ para todo $J$, y $\PP(A)=(1-p)^{k}$.
  \item Igualando 3 y 4: $\E\bigl[\binom{n-C}{k}/\binom nk\bigr]=(1-p)^{k}$, y por linealidad $\E[\widehat\pi_k]=1-(1-p)^{k}$.
\end{enumerate}
\resultado \textbf{$\widehat\pi_k$ es un estimador insesgado de pass@$k$.} $\blacksquare$
\comprobacion para $k=1$, $\widehat\pi_1=1-(n-C)/n=C/n$, la proporción muestral, que es insesgada para $p$. Numéricamente, con $n=10$, $k=3$ y $p=0{,}3$, la esperanza calculada sumando sobre $C=0,\dots,10$ vale $0{,}657=1-0{,}7^{3}$.
\end{solucion}

% ---------------------------------------------------------------------
\section{Tema 3}

\begin{solucion}{eje:t3-vf}
\planteamiento contrastar con el \cref{teo:t3-tcl}, el \cref{teo:t3-suma-normales}, el \cref{cor:t3-binomial} y la \cref{def:t3-chi-t}.
\pasos
\begin{enumerate}[label=(\alph*)]
  \item \textbf{Falsa.} Los datos siguen siendo exponenciales (asimétricos); lo que es aproximadamente normal es la distribución de $\Xbar$.
  \item \textbf{Verdadera}, por el \cref{teo:t3-suma-normales}.
  \item \textbf{Verdadera}: $\ee(\Xbar)=\sigma/\sqrt{100}=\sigma/10$.
  \item \textbf{Verdadera}: por ejemplo, $t_{10;0{,}975}=2{,}228>1{,}960=z_{0{,}975}$, de modo que $\PP(T_{10}>1{,}96)>0{,}025$.
  \item \textbf{Falsa}: $np=0{,}4$, muy lejos de $np\ge10$; la distribución es muy asimétrica (de hecho, $\PP(X=0)=0{,}98^{20}\approx0{,}668$).
\end{enumerate}
\resultado \textbf{(a) F; (b) V; (c) V; (d) V; (e) F.}
\end{solucion}

\begin{solucion}{eje:t3-tabla}
\planteamiento leer la \cref{tab:phi} y usar $\Phi(-z)=1-\Phi(z)$.
\pasos
\begin{enumerate}[label=(\alph*)]
  \item Fila $1{,}2$, columna $0{,}03$: $\Phi(1{,}23)=0{,}8907$.
  \item $\PP(Z>-0{,}5)=1-\Phi(-0{,}5)=\Phi(0{,}5)=0{,}6915$.
  \item $\Phi(2)-\Phi(-1)=0{,}9772-(1-0{,}8413)=0{,}8185$.
  \item $z_{0{,}95}=1{,}6449$ y $z_{0{,}99}=2{,}3263$ (\cref{tab:zq}).
  \item $\PP(\abs Z\le c)=2\Phi(c)-1=0{,}90\iff\Phi(c)=0{,}95\iff c=z_{0{,}95}=1{,}6449$.
\end{enumerate}
\resultado \textbf{(a) $0{,}8907$; (b) $0{,}6915$; (c) $0{,}8185$; (d) $1{,}6449$ y $2{,}3263$; (e) $c\approx1{,}645$.}
\comprobacion en (e), cada cola tiene probabilidad $0{,}05$, y por eso aparece el cuantil $0{,}95$ y no el $0{,}90$.
\end{solucion}

\begin{solucion}{eje:t3-normal-basico}
\datos $X\sim\Nor(100,15^{2})$.
\planteamiento tipificar (\cref{prop:t3-tipificacion}).
\pasos
\begin{enumerate}[label=\arabic*.]
  \item $\PP(X>130)=1-\Phi\bigl(\frac{130-100}{15}\bigr)=1-\Phi(2)=1-0{,}9772=0{,}0228$.
  \item $\PP(85<X<115)=\Phi(1)-\Phi(-1)=2\Phi(1)-1=2\cdot0{,}8413-1=0{,}6826$.
  \item Percentil $90$: $x=100+15\,z_{0{,}90}=100+15\cdot1{,}2816=119{,}22$.
\end{enumerate}
\resultado \textbf{$0{,}0228$; $0{,}683$; $119{,}2$.}
\comprobacion los dos primeros resultados reproducen la regla empírica ($68\,\%$ a una desviación típica; $2{,}3\,\%$ por encima de dos).
\end{solucion}

\begin{solucion}{eje:t3-openstax-media}
\datos $\mu=90$, $\sigma=15$, $n=25$.
\planteamiento por el teorema central del límite, $\Xbar\approx\Nor(90,\,3^{2})$, pues $\ee=15/\sqrt{25}=3$.
\pasos
\begin{enumerate}[label=\arabic*.]
  \item $z_1=(85-90)/3\approx-1{,}67$; $z_2=(92-90)/3\approx0{,}67$.
  \item $\PP(85<\Xbar<92)\approx\Phi(0{,}67)-\Phi(-1{,}67)=0{,}7486-0{,}0475=0{,}7011$.
  \item Dos errores estándar por encima: $90+2\cdot3=96$.
\end{enumerate}
\resultado \textbf{(a) $\approx0{,}70$; (b) $96$.}
\comprobacion sin redondear $z$ a dos decimales se obtiene $0{,}6997$, que es el valor de la fuente.
\end{solucion}

\begin{solucion}{eje:t3-openstax-suma}
\datos $\mu=90$, $\sigma=15$, $n=80$.
\planteamiento $S_n\approx\Nor(n\mu,n\sigma^{2})$ (\cref{obs:t3-uso}).
\pasos
\begin{enumerate}[label=\arabic*.]
  \item $\E S_n=80\cdot90=7200$; $\sqrt{\Var S_n}=\sqrt{80}\cdot15\approx134{,}16$.
  \item $z=(7500-7200)/134{,}16\approx2{,}24$; $\PP(S_n>7500)\approx1-\Phi(2{,}24)=1-0{,}9875=0{,}0125$.
  \item $7200+1{,}5\cdot134{,}16\approx7401{,}2$.
\end{enumerate}
\resultado \textbf{(a) $\approx0{,}0125$; (b) $\approx7401$.}
\comprobacion con $z=2{,}236$ sin redondear se obtiene $0{,}0127$, el valor de la fuente.
\end{solucion}

\begin{solucion}{eje:t3-elecciones}
\datos $n=400$; $p_A=0{,}5$; $p_{\text{menores}}=0{,}3$.
\planteamiento $\widehat p\approx\Nor\bigl(p,p(1-p)/n\bigr)$ (\cref{cor:t3-binomial}).
\pasos
\begin{enumerate}[label=\arabic*.]
  \item (a) $\ee=\sqrt{0{,}25/400}=0{,}025$; $z=(0{,}525-0{,}5)/0{,}025=1$; $\PP(\widehat p\ge0{,}525)\approx1-\Phi(1)=0{,}1587$.
  \item (b) $\ee=\sqrt{0{,}21/400}\approx0{,}02291$; $z=(0{,}31-0{,}3)/0{,}02291\approx0{,}44$; $\PP(\widehat p<0{,}31)\approx\Phi(0{,}44)=0{,}6700$.
\end{enumerate}
\resultado \textbf{(a) $\approx0{,}16$; (b) $\approx0{,}67$.}
\comprobacion en (a), con corrección por continuidad, $\PP(X\ge210)\approx1-\Phi\bigl((209{,}5-200)/10\bigr)=1-\Phi(0{,}95)=0{,}1711$, que coincide con el valor binomial exacto ($0{,}1711$).
\end{solucion}

\begin{solucion}{eje:t3-redondeo}
\datos errores $E_i$ uniformes en $\{-2,-1,0,1,2\}$, independientes; $n=1000$.
\planteamiento calcular media y varianza de cada error y aplicar el teorema central del límite a la suma $S$.
\pasos
\begin{enumerate}[label=\arabic*.]
  \item $\E E_i=0$; $\Var E_i=\E E_i^{2}=(4+1+0+1+4)/5=2$.
  \item $\E S=0$; $\Var S=1000\cdot2=2000$; desviación típica $\sqrt{2000}\approx44{,}72$.
  \item $\PP(\abs S>100)\approx2\bigl(1-\Phi(100/44{,}72)\bigr)=2\bigl(1-\Phi(2{,}24)\bigr)=2\cdot0{,}0125=0{,}025$.
\end{enumerate}
\resultado \textbf{$p\approx0{,}025$.}
\comprobacion como $S$ es entera, $\abs S>100$ equivale a $\abs S\ge101$, y con corrección por continuidad se obtiene $2\bigl(1-\Phi(100{,}5/44{,}72)\bigr)\approx0{,}0246$; la diferencia es pequeña.
\end{solucion}

\begin{solucion}{eje:t3-binomial}
\datos $np=50$, $\sqrt{np(1-p)}=5$.
\planteamiento \cref{cor:t3-binomial} y \cref{def:t3-cc}.
\pasos
\begin{enumerate}[label=\arabic*.]
  \item $\PP(X=50)=\PP(49{,}5\le Y\le50{,}5)\approx\Phi(0{,}1)-\Phi(-0{,}1)=2\cdot0{,}5398-1=0{,}0796$.
  \item Con corrección: $\PP(39{,}5\le Y\le60{,}5)\approx\Phi(2{,}1)-\Phi(-2{,}1)=2\cdot0{,}9821-1=0{,}9642$.
  \item Sin corrección: $\PP(40\le Y\le60)\approx2\Phi(2)-1=2\cdot0{,}9772-1=0{,}9544$.
\end{enumerate}
\resultado \textbf{(a) $0{,}0796$; (b) $0{,}9642$ con corrección y $0{,}9544$ sin ella.}
\comprobacion los valores exactos son $0{,}0796$ y $0{,}9648$: con corrección el error es de $0{,}0006$ y sin ella de $0{,}0104$. Sin corrección, $\PP(X=50)$ se aproximaría por $0$.
\end{solucion}

\begin{solucion}{eje:t3-benchmark}
\datos $p=0{,}8$; $n=1000$.
\planteamiento $\widehat p\approx\Nor\bigl(0{,}8,\ 0{,}16/1000\bigr)$.
\pasos
\begin{enumerate}[label=\arabic*.]
  \item $\ee=\sqrt{0{,}8\cdot0{,}2/1000}=\sqrt{0{,}000\,16}\approx0{,}01265$.
  \item $z=(0{,}82-0{,}8)/0{,}01265\approx1{,}58$; $\PP(\widehat p\ge0{,}82)\approx1-\Phi(1{,}58)=1-0{,}9429=0{,}0571$.
  \item Intervalo: $0{,}8\pm z_{0{,}975}\,\ee=0{,}8\pm1{,}96\cdot0{,}01265=0{,}8\pm0{,}0248$, es decir, $[0{,}775;\allowbreak\ 0{,}825]$.
\end{enumerate}
\resultado \textbf{(a) $0{,}0126$; (b) $\approx0{,}057$; (c) $[0{,}775;\allowbreak\,0{,}825]$.}
\comprobacion un modelo con exactitud real del $80\,\%$ obtiene un $82\,\%$ o más en casi un $6\,\%$ de las evaluaciones: una ventaja de dos puntos en un benchmark de $1000$ preguntas no es concluyente por sí sola.
\end{solucion}

\begin{solucion}{eje:t3-tamano}
\datos probabilidad $0{,}95$; tolerancia $0{,}1\sigma$; en (b), tolerancia $0{,}05$ y $\sigma=\sqrt{p(1-p)}\le1/2$.
\planteamiento $\PP(\abs{\Xbar-\mu}<c)\approx2\Phi\bigl(c\sqrt n/\sigma\bigr)-1\ge0{,}95\iff c\sqrt n/\sigma\ge z_{0{,}975}=1{,}96$.
\pasos
\begin{enumerate}[label=\arabic*.]
  \item (a) $c=0{,}1\sigma$: $0{,}1\sqrt n\ge1{,}96\iff n\ge(19{,}6)^{2}=384{,}16$, luego $n=385$.
  \item (b) $c=0{,}05$ y $\sigma\le0{,}5$: basta $0{,}05\sqrt n/0{,}5\ge1{,}96$, que es la misma condición que en (a): $n=385$.
\end{enumerate}
\resultado \textbf{$n=385$ en ambos casos}, frente a $n=2000$ con Chebyshev (\cref{eje:t2-chebyshev-n}).
\comprobacion con $n=385$ y $p=1/2$: $\ee=0{,}5/\sqrt{385}\approx0{,}02548$ y $1{,}96\cdot0{,}02548\approx0{,}0499<0{,}05$.
\end{solucion}

\begin{solucion}{eje:t3-t-cuantiles}
\planteamiento leer las \cref{tab:t,tab:chi2} y usar la simetría de la $t$.
\pasos
\begin{enumerate}[label=\arabic*.]
  \item $t_{9;0{,}975}=2{,}262$, $t_{29;0{,}975}=2{,}045$, $t_{120;0{,}975}=1{,}980$, frente a $z_{0{,}975}=1{,}960$.
  \item $t_{15;0{,}975}=2{,}131$, luego $\PP(T>2{,}131)=0{,}025$.
  \item $t_{15;0{,}95}=1{,}753$, luego $\PP(\abs T\le1{,}753)=1-2\cdot0{,}05=0{,}90$.
  \item $\PP(V>c)=0{,}05\iff c=\chi^{2}_{10;0{,}95}=18{,}307$.
\end{enumerate}
\resultado \textbf{(a) $2{,}262$; $2{,}045$; $1{,}980$ (decrecen hacia $1{,}960$); (b) $0{,}025$ y $0{,}90$; (c) $c=18{,}307$.}
\comprobacion la $\chi^{2}_{10}$ tiene media $10$ y desviación típica $\sqrt{20}\approx4{,}47$; $18{,}3$ está unas $1{,}9$ desviaciones por encima, razonable para un cuantil $0{,}95$ de una distribución asimétrica a la derecha.
\end{solucion}

\begin{solucion}{eje:t3-lineal-normal}
\planteamiento calcular la función de distribución de $Y$, derivarla y reconocer la densidad normal.
\pasos
\begin{enumerate}[label=\arabic*.]
  \item Caso $a>0$: $F_Y(y)=\PP(aX+b\le y)=\PP\bigl(X\le\frac{y-b}{a}\bigr)=F_X\bigl(\frac{y-b}{a}\bigr)$. Derivando (regla de la cadena), $f_Y(y)=\frac1af_X\bigl(\frac{y-b}{a}\bigr)$.
  \item Caso $a<0$: $F_Y(y)=\PP\bigl(X\ge\frac{y-b}{a}\bigr)=1-F_X\bigl(\frac{y-b}{a}\bigr)$, y $f_Y(y)=-\frac1af_X\bigl(\frac{y-b}{a}\bigr)$. En ambos casos, $f_Y(y)=\frac{1}{\abs a}f_X\bigl(\frac{y-b}{a}\bigr)$.
  \item Sustituyendo la densidad normal y usando $\frac{y-b}{a}-\mu=\frac{y-(a\mu+b)}{a}$:
  \[
  f_Y(y)=\frac{1}{\abs a\sigma\sqrt{2\pi}}\exp\Bigl(-\frac{(y-(a\mu+b))^{2}}{2a^{2}\sigma^{2}}\Bigr),
  \]
  que es la densidad $\Nor(a\mu+b,a^{2}\sigma^{2})$, con desviación típica $\abs a\sigma$.
  \item Con $a=1/\sigma$ y $b=-\mu/\sigma$ se obtiene $Z=(X-\mu)/\sigma\sim\Nor(0,1)$ y $\PP(X\le x)=\PP\bigl(Z\le\frac{x-\mu}{\sigma}\bigr)=\Phi\bigl(\frac{x-\mu}{\sigma}\bigr)$.
\end{enumerate}
\resultado \textbf{$aX+b\sim\Nor(a\mu+b,a^{2}\sigma^{2})$}, y la tipificación lleva cualquier normal a la estándar. $\blacksquare$
\comprobacion con $a=-1$, $b=0$, $\mu=0$, $\sigma=1$: $-Z\sim\Nor(0,1)$, que expresa la simetría de la normal estándar.
\end{solucion}

\begin{solucion}{eje:t3-chi1}
\planteamiento para $v>0$, $Z^{2}\le v\iff-\sqrt v\le Z\le\sqrt v$.
\pasos
\begin{enumerate}[label=\arabic*.]
  \item $\PP(V\le v)=\Phi(\sqrt v)-\Phi(-\sqrt v)=\Phi(\sqrt v)-\bigl(1-\Phi(\sqrt v)\bigr)=2\Phi(\sqrt v)-1$.
  \item Sea $\varphi(z)=\frac{1}{\sqrt{2\pi}}e^{-z^{2}/2}=\Phi'(z)$. Por la regla de la cadena,
  \[
  f_V(v)=2\varphi(\sqrt v)\cdot\frac{1}{2\sqrt v}=\frac{\varphi(\sqrt v)}{\sqrt v}=\frac{1}{\sqrt{2\pi v}}e^{-v/2},\qquad v>0.
  \]
  \item Para $v\le0$, $\PP(V\le v)=\PP(Z^{2}\le v)\le\PP(Z=0)=0$, y la densidad es nula.
\end{enumerate}
\resultado \textbf{$f_V(v)=(2\pi v)^{-1/2}e^{-v/2}$ para $v>0$.} $\blacksquare$
\comprobacion $\PP(V\le3{,}841)=2\Phi(\sqrt{3{,}841})-1=2\Phi(1{,}96)-1=0{,}95$, de acuerdo con $\chi^{2}_{1;0{,}95}=3{,}841$ en la \cref{tab:chi2}.
\end{solucion}

\begin{solucion}{eje:t3-chi-momentos}
\planteamiento $\varphi'(z)=-z\varphi(z)$; integrar por partes $\int z^{m}\cdot z\varphi(z)\,dz$; después, linealidad e independencia.
\pasos
\begin{enumerate}[label=\arabic*.]
  \item Para $m\ge0$ entero, con $u=z^{m+1}$ y $dv=z\varphi(z)\,dz$, $v=-\varphi(z)$:
  \[
  \begin{aligned}
  \E[Z^{m+2}]&=\int_{-\infty}^{\infty}z^{m+1}\,z\varphi(z)\,dz=\Bigl[-z^{m+1}\varphi(z)\Bigr]_{-\infty}^{\infty}+(m+1)\int_{-\infty}^{\infty}z^{m}\varphi(z)\,dz\\
  &=(m+1)\E[Z^{m}],
  \end{aligned}
  \]
  porque $z^{m+1}e^{-z^{2}/2}\to0$ cuando $\abs z\to\infty$.
  \item $m=0$: $\E[Z^{2}]=1\cdot\E[1]=1$. $m=2$: $\E[Z^{4}]=3\E[Z^{2}]=3$.
  \item $\Var(Z^{2})=\E[Z^{4}]-(\E Z^{2})^{2}=3-1=2$.
  \item $V=\sum_{i=1}^{k}Z_i^{2}$ con $Z_i$ \iid: $\E V=k\cdot1=k$ y, por independencia, $\Var V=k\cdot2=2k$.
\end{enumerate}
\resultado \textbf{$\E[Z^{4}]=3$; $\E[\chi^{2}_k]=k$ y $\Var(\chi^{2}_k)=2k$.} $\blacksquare$
\comprobacion para $k=10$, media $10$ y varianza $20$, coherente con el \cref{eje:t3-t-cuantiles}.
\end{solucion}

\begin{solucion}{eje:t3-t-simetrica}
\planteamiento escribir $T=Z/\sqrt{V/k}$ con $Z\sim\Nor(0,1)$ y $V\sim\chi^{2}_k$ independientes, y observar que $-T$ es la misma función aplicada a $(-Z,V)$.
\pasos
\begin{enumerate}[label=\arabic*.]
  \item $-Z\sim\Nor(0,1)$ (\cref{eje:t3-lineal-normal} con $a=-1$) y $-Z$ es independiente de $V$. Luego el par $(-Z,V)$ tiene la misma distribución conjunta que $(Z,V)$.
  \item Como $-T=(-Z)/\sqrt{V/k}$ es la misma función del par $(-Z,V)$ que $T$ del par $(Z,V)$, $-T$ y $T$ tienen la misma distribución.
  \item Sea $F$ la función de distribución de $T$, continua. Para todo $t$: $F(-t)=\PP(T\le-t)=\PP(-T\ge t)=\PP(T\ge t)=1-F(t)$.
  \item Con $t=t_{k;q}$: $F(-t_{k;q})=1-q$, es decir, $-t_{k;q}=t_{k;1-q}$.
  \item Con $c=t_{k;1-\alpha/2}$: $\PP(\abs T\le c)=F(c)-F(-c)=2F(c)-1=2(1-\alpha/2)-1=1-\alpha$.
\end{enumerate}
\resultado \textbf{$T\stackrel{d}{=}-T$, $t_{k;1-q}=-t_{k;q}$ y $\PP(\abs T\le t_{k;1-\alpha/2})=1-\alpha$.} $\blacksquare$
\comprobacion el mismo argumento con $V=k$ constante da la simetría de la normal estándar, cuyo caso límite es la $t_k$ cuando $k\to\infty$.
\end{solucion}

% ---------------------------------------------------------------------
\section{Tema 4}

\begin{solucion}{eje:t4-interpretacion}
\planteamiento aplicar la \cref{obs:t4-interpretacion}: la confianza se refiere al procedimiento, no a la probabilidad de un valor concreto del parámetro.
\pasos
\begin{enumerate}[label=(\roman*)]
  \item \textbf{No se deduce.} En el enfoque frecuentista la mejora real es fija; el $95\,\%$ es la cobertura del procedimiento.
  \item \textbf{No se deduce}, por la misma razón: el intervalo no asigna probabilidades a hipótesis.
  \item \textbf{Se deduce}: el centro, $15$, es la estimación puntual y el intervalo indica su precisión.
  \item \textbf{No se deduce}: «probablemente falsa» es una afirmación de probabilidad sobre una hipótesis.
  \item \textbf{Se deduce}: es la definición de cobertura (\cref{def:t4-ic}).
  \item \textbf{Se deduce}: el valor $0$ queda fuera del intervalo, y un intervalo de nivel $95\,\%$ contiene exactamente los valores no rechazados por el contraste bilateral al $5\,\%$ (\cref{eje:t5-dualidad}).
\end{enumerate}
\resultado \textbf{Se deducen (iii), (v) y (vi); no se deducen (i), (ii) y (iv).}
\end{solucion}

\begin{solucion}{eje:t4-momentos-basico}
\planteamiento \cref{ej:t4-momentos}: igualar media (y segundo momento) muestrales y poblacionales.
\pasos
\begin{enumerate}[label=\arabic*.]
  \item (a) $\xbar=(0{,}5+1{,}2+0{,}3+2{,}0)/4=4{,}0/4=1{,}0$; $\widehat\lambda=1/\xbar=1{,}0$ peticiones por segundo.
  \item (b) $\widehat\mu=\xbar=22/5=4{,}4$. Desviaciones: $-2{,}4;\allowbreak\,-0{,}4;\allowbreak\,-0{,}4;\allowbreak\,0{,}6;\allowbreak\,2{,}6$; cuadrados: $5{,}76;\allowbreak\,0{,}16;\allowbreak\,0{,}16;\allowbreak\,0{,}36;\allowbreak\,6{,}76$; suma $13{,}2$.
  \item $\widehat\sigma^{2}=13{,}2/5=2{,}64$, mientras que $s^{2}=13{,}2/4=3{,}3$.
\end{enumerate}
\resultado \textbf{(a) $\widehat\lambda=1$; (b) $\widehat\mu=4{,}4$ y $\widehat\sigma^{2}=2{,}64$, menor que $s^{2}=3{,}3$.}
\comprobacion $\frac1n\sum x_i^{2}-\xbar^{2}=(4+16+16+25+49)/5-19{,}36=22-19{,}36=2{,}64$.
\end{solucion}

\begin{solucion}{eje:t4-geometrica}
\datos $f(x\mid p)=p(1-p)^{x-1}$, $\E X=1/p$; datos $2,1,4,3,1,1$.
\planteamiento (a) $1/p=\xbar$; (b) derivar la log-verosimilitud.
\pasos
\begin{enumerate}[label=\arabic*.]
  \item Momentos: $1/\widehat p=\xbar$, luego $\widehat p=1/\xbar$.
  \item Verosimilitud: $\Lik(p)=\prod_ip(1-p)^{x_i-1}=p^{n}(1-p)^{\sum x_i-n}$; $\ell(p)=n\ln p+\bigl(\sum x_i-n\bigr)\ln(1-p)$.
  \item $\ell'(p)=\dfrac np-\dfrac{\sum x_i-n}{1-p}=0\iff n(1-p)=\bigl(\sum x_i-n\bigr)p\iff n=p\sum x_i\iff p=\dfrac{n}{\sum x_i}=\dfrac1{\xbar}$.
  \item $\ell''(p)=-n/p^{2}-\bigl(\sum x_i-n\bigr)/(1-p)^{2}<0$: máximo global.
  \item Datos: $\sum x_i=12$, $n=6$, $\xbar=2$, $\widehat p=0{,}5$.
\end{enumerate}
\resultado \textbf{Ambos métodos dan $\widehat p=1/\xbar$; con los datos, $\widehat p=0{,}5$.}
\comprobacion un agente que resuelve cada intento con probabilidad $0{,}5$ necesita en media $2$ intentos, que es la media observada.
\end{solucion}

\begin{solucion}{eje:t4-gamma}
\datos $f(x\mid a)=a^{5}x^{4}e^{-ax}/24$; datos $2{,}0;\allowbreak\,3{,}0;\allowbreak\,2{,}5;\allowbreak\,4{,}5$.
\planteamiento escribir $\ell(a)$ con $S=\sum x_i$ y $P=\prod x_i$ y resolver $\ell'(a)=0$.
\pasos
\begin{enumerate}[label=\arabic*.]
  \item $\Lik(a)=\dfrac{a^{5n}P^{4}e^{-aS}}{24^{n}}$; $\ell(a)=5n\ln a+4\ln P-aS-n\ln24$.
  \item $\ell'(a)=5n/a-S=0\iff a=5n/S$.
  \item $\ell''(a)=-5n/a^{2}<0$: máximo global en $(0,\infty)$.
  \item Datos: $n=4$, $S=12$; $\widehat a=20/12=5/3\approx1{,}667$.
\end{enumerate}
\resultado \textbf{$\widehat a=5n/\sum x_i=5/\xbar$; con los datos, $\widehat a=5/3$.}
\comprobacion la media de esta distribución es $5/a$ (es una gamma de forma $5$), luego el EMV coincide con el estimador por momentos: $5/\widehat a=\xbar=3$.
\end{solucion}

\begin{solucion}{eje:t4-uniforme-ab}
\datos $1{,}2;\allowbreak\,2{,}1;\allowbreak\,1{,}3;\allowbreak\,10{,}5;\allowbreak\,5$; $n=5$.
\planteamiento (a) $\Lik(a,b)=(b-a)^{-n}$ si $a\le\min x_i$ y $b\ge\max x_i$, y $0$ en otro caso; (b) igualar media y varianza (con denominador $n$).
\pasos
\begin{enumerate}[label=\arabic*.]
  \item (a) $\Lik$ es positiva solo si el intervalo $[a,b]$ contiene todos los datos y, en ese caso, es mayor cuanto menor es $b-a$. El intervalo más corto que contiene los datos es $[\min x_i,\max x_i]$: $\widehat a=1{,}2$ y $\widehat b=10{,}5$.
  \item (b) $\xbar=20{,}1/5=4{,}02$; $\frac1n\sum x_i^{2}=(1{,}44+4{,}41+1{,}69+110{,}25+25)/5=142{,}79/5=28{,}558$; $\widehat\sigma^{2}=28{,}558-4{,}02^{2}=12{,}3976$; $\widehat\sigma\approx3{,}521$.
  \item De $(b-a)/2=\sqrt3\,\widehat\sigma\approx6{,}099$ y $(a+b)/2=4{,}02$: $\widehat a\approx-2{,}08$ y $\widehat b\approx10{,}12$.
  \item (c) $\widehat b\approx10{,}12<10{,}5$: según el modelo estimado, el dato $10{,}5$ sería imposible.
\end{enumerate}
\resultado \textbf{EMV: $\widehat a=1{,}2$, $\widehat b=10{,}5$. Momentos: $\widehat a\approx-2{,}08$, $\widehat b\approx10{,}12$, incompatibles con los datos.}
\comprobacion la media del modelo con los estimadores de momentos es $(-2{,}08+10{,}12)/2=4{,}02=\xbar$, como debe ser.
\end{solucion}

\begin{solucion}{eje:t4-poisson-datos}
\datos $2,0,3,1,4,2,1,3$; $n=8$.
\planteamiento $\widehat\lambda=\xbar$ (\cref{eje:t4-emv-poisson}); intervalo $\widehat\lambda\pm1{,}96\sqrt{\widehat\lambda/n}$.
\pasos
\begin{enumerate}[label=\arabic*.]
  \item $\sum x_i=16$; $\widehat\lambda=2$.
  \item $\widehat{\ee}=\sqrt{2/8}=0{,}5$; margen $1{,}96\cdot0{,}5=0{,}98$.
\end{enumerate}
\resultado \textbf{$\widehat\lambda=2$; intervalo aproximado $[1{,}02;\allowbreak\ 2{,}98]$.}
\comprobacion con $n=8$ la aproximación normal es tosca; el intervalo es solo orientativo, pero su amplitud ya muestra que $\lambda$ está poco determinado.
\end{solucion}

\begin{solucion}{eje:t4-soda}
\datos $n=9$; $\sum x_i=3204$; $\xbar=356$.
\planteamiento (a) \cref{teo:t4-ic-z}; (b) \cref{teo:t4-ic-t} con $8$ grados de libertad.
\pasos
\begin{enumerate}[label=\arabic*.]
  \item $\sigma/\sqrt n=3/3=1$. Al $95\,\%$: $356\pm1{,}96$, es decir, $[354{,}04;\allowbreak\ 357{,}96]$. Al $98\,\%$ ($z_{0{,}99}=2{,}326$): $[353{,}67;\allowbreak\ 358{,}33]$.
  \item Desviaciones respecto de $356$: $-4,-5,5,-3,-4,2,4,2,3$; suma de cuadrados $16+25+25+9+16+4+16+4+9=124$; $s^{2}=124/8=15{,}5$; $s\approx3{,}937$; $s/\sqrt n\approx1{,}312$.
  \item Al $95\,\%$: $t_{8;0{,}975}=2{,}306$, margen $3{,}026$, intervalo $[352{,}97;\allowbreak\ 359{,}03]$.
  \item Al $98\,\%$: $t_{8;0{,}99}=2{,}896$, margen $3{,}801$, intervalo $[352{,}20;\allowbreak\ 359{,}80]$.
\end{enumerate}
\resultado \textbf{$\sigma$ conocida: $[354{,}04;\allowbreak\,357{,}96]$ y $[353{,}67;\allowbreak\,358{,}33]$. $\sigma$ desconocida: $[352{,}97;\allowbreak\,359{,}03]$ y $[352{,}20;\allowbreak\,359{,}80]$.}
\comprobacion los intervalos $t$ son más anchos por dos motivos: $t_{8}>z$ y, en esta muestra, $s\approx3{,}94>\sigma=3$. Al aumentar el nivel, todos los intervalos se ensanchan.
\end{solucion}

\begin{solucion}{eje:t4-alturas}
\datos $n=20$; $\xbar=69{,}55$; $s^{2}=14{,}26$; en (b), $\sigma=3{,}77$.
\planteamiento \cref{teo:t4-ic-t,teo:t4-ic-z} y \cref{cor:t4-tamano} con $E=0{,}5$ (la mitad de la amplitud).
\pasos
\begin{enumerate}[label=\arabic*.]
  \item (a) $s/\sqrt n=\sqrt{14{,}26/20}=\sqrt{0{,}713}\approx0{,}8444$; $t_{19;0{,}95}=1{,}729$; margen $1{,}460$; intervalo $[68{,}09;\allowbreak\ 71{,}01]$.
  \item (b) $\sigma/\sqrt n=3{,}77/\sqrt{20}\approx0{,}8430$; $z_{0{,}95}=1{,}645$; margen $1{,}387$; intervalo $[68{,}16;\allowbreak\ 70{,}94]$.
  \item (c) $n\ge(1{,}645\cdot3{,}77/0{,}5)^{2}=(12{,}40)^{2}\approx153{,}8$, luego $n=154$.
\end{enumerate}
\resultado \textbf{(a) $[68{,}09;\allowbreak\,71{,}01]$; (b) $[68{,}16;\allowbreak\,70{,}94]$; (c) $154$ personas.}
\comprobacion con $n=154$ la amplitud es $2\cdot1{,}645\cdot3{,}77/\sqrt{154}\approx0{,}9995<1$.
\end{solucion}

\begin{solucion}{eje:t4-openstax-z}
\datos (a) $\sigma=3$, $n=36$, $\xbar=68$, nivel $90\,\%$; (b) $\sigma=15$, $E=2$, nivel $95\,\%$.
\planteamiento \cref{teo:t4-ic-z} y \cref{cor:t4-tamano}.
\pasos
\begin{enumerate}[label=\arabic*.]
  \item (a) Margen $1{,}645\cdot3/6=0{,}8225$; intervalo $[67{,}18;\allowbreak\ 68{,}82]$.
  \item (b) $n\ge(1{,}96\cdot15/2)^{2}=14{,}7^{2}=216{,}09$, luego $n=217$.
\end{enumerate}
\resultado \textbf{(a) $[67{,}18;\allowbreak\,68{,}82]$; (b) $n=217$.}
\comprobacion se redondea hacia arriba: con $n=216$ el margen sería $1{,}96\cdot15/\sqrt{216}\approx2{,}0004>2$.
\end{solucion}

\begin{solucion}{eje:t4-proporcion}
\datos (a) $n=500$, $x=421$; (b) $E=0{,}03$, nivel $90\,\%$.
\planteamiento \cref{def:t4-wald} y \cref{cor:t4-tamano-p}.
\pasos
\begin{enumerate}[label=\arabic*.]
  \item (a) $\widehat p=421/500=0{,}842$; $\widehat{\ee}=\sqrt{0{,}842\cdot0{,}158/500}\approx0{,}01631$; margen $1{,}96\cdot0{,}01631\approx0{,}0320$; intervalo $[0{,}810;\allowbreak\ 0{,}874]$.
  \item (b) $n\ge z_{0{,}95}^{2}/(4E^{2})=1{,}645^{2}/(4\cdot0{,}0009)=2{,}706/0{,}0036\approx751{,}7$, luego $n=752$.
\end{enumerate}
\resultado \textbf{(a) $[0{,}810;\allowbreak\,0{,}874]$; (b) $752$ personas.}
\comprobacion en (a), $n\widehat p=421$ y $n(1-\widehat p)=79$ son grandes, y la aproximación normal es adecuada.
\end{solucion}

\begin{solucion}{eje:t4-margen}
\planteamiento margen conservador $z/(2\sqrt n)=E\iff n=\bigl(z/(2E)\bigr)^{2}$.
\pasos
\begin{enumerate}[label=\arabic*.]
  \item (a) $z=1{,}96$: $n=(1{,}96/0{,}02)^{2}=98^{2}=9604$. Con $z\approx2$: $n=100^{2}=10\,000$.
  \item (b) $z_{0{,}90}=1{,}2816$: $n=(1{,}2816/0{,}02)^{2}=64{,}08^{2}\approx4106$.
  \item (c) Con $z\approx2$: $2/(2\sqrt n)=0{,}04\iff\sqrt n=25\iff n=625$ (con $1{,}96$ saldría unas $600$).
\end{enumerate}
\resultado \textbf{(a) $9604$ (o $10\,000$ con $z\approx2$); (b) unas $4106$; (c) unas $600$ personas.}
\comprobacion la regla práctica «margen $\approx1/\sqrt n$ al $95\,\%$» da $1/\sqrt{625}=0{,}04$: un benchmark de $10\,000$ preguntas tiene un margen de un punto porcentual.
\end{solucion}

\begin{solucion}{eje:t4-emv-poisson}
\planteamiento escribir $\ell(\lambda)$, resolver $\ell'(\lambda)=0$ y estudiar la concavidad.
\pasos
\begin{enumerate}[label=\arabic*.]
  \item $f(x\mid\lambda)=e^{-\lambda}\lambda^{x}/x!$, luego $\ell(\lambda)=\Bigl(\sum_ix_i\Bigr)\ln\lambda-n\lambda-\sum_i\ln(x_i!)$ para $\lambda>0$.
  \item $\ell'(\lambda)=\dfrac{\sum_ix_i}{\lambda}-n=0\iff\lambda=\dfrac{\sum_ix_i}{n}=\xbar$.
  \item $\ell''(\lambda)=-\sum_ix_i/\lambda^{2}<0$ para todo $\lambda>0$ (pues $\sum x_i>0$): $\ell$ es estrictamente cóncava en $(0,\infty)$ y su único punto crítico es el máximo global.
  \item Por el \cref{ej:t4-momentos}, el estimador por momentos también es $\Xbar$.
\end{enumerate}
\resultado \textbf{El EMV es $\widehat\lambda=\xbar$ y coincide con el de momentos.} $\blacksquare$
\comprobacion con los datos del \cref{eje:t4-poisson-datos}, $\ell'(2)=16/2-8=0$.
\end{solucion}

\begin{solucion}{eje:t4-emv-expo}
\planteamiento derivar la log-verosimilitud en cada parametrización.
\pasos
\begin{enumerate}[label=\arabic*.]
  \item (a) $\ell(\lambda)=n\ln\lambda-\lambda\sum_ix_i$; $\ell'(\lambda)=n/\lambda-\sum_ix_i=0\iff\lambda=n/\sum_ix_i=1/\xbar$; $\ell''(\lambda)=-n/\lambda^{2}<0$: máximo global.
  \item (b) Con $\beta=1/\lambda$, $f(x\mid\beta)=\beta^{-1}e^{-x/\beta}$ y $\ell(\beta)=-n\ln\beta-\sum_ix_i/\beta$.
  \item $\ell'(\beta)=-\dfrac n\beta+\dfrac{\sum_ix_i}{\beta^{2}}=\dfrac{\sum_ix_i-n\beta}{\beta^{2}}$, positiva para $\beta<\xbar$ y negativa para $\beta>\xbar$: máximo global en $\widehat\beta=\xbar$.
  \item $\widehat\beta=\xbar=1/\widehat\lambda$, como predice la invariancia con $g(\lambda)=1/\lambda$.
\end{enumerate}
\resultado \textbf{$\widehat\lambda=1/\xbar$ y $\widehat\beta=\xbar$.} $\blacksquare$
\comprobacion con los datos del \cref{eje:t4-momentos-basico}(a), $\widehat\lambda=1$ y $\widehat\beta=1$ segundo.
\end{solucion}

\begin{solucion}{eje:t4-laplace}
\planteamiento la log-verosimilitud es, salvo una constante, $-\sum_i\abs{x_i-\mu}$.
\pasos
\begin{enumerate}[label=\arabic*.]
  \item $\ell(\mu)=\sum_i\ln\bigl(\tfrac12e^{-\abs{x_i-\mu}}\bigr)=-n\ln2-\sum_i\abs{x_i-\mu}$.
  \item Maximizar $\ell$ equivale a minimizar $g(\mu)=\sum_i\abs{x_i-\mu}$, que alcanza su mínimo en la mediana (\cref{eje:t1-mediana-min}).
  \item En el modelo normal con $\sigma$ fija, $\ell(\mu)=c-\frac{1}{2\sigma^{2}}\sum_i(x_i-\mu)^{2}$, que se maximiza en $\xbar$ (\cref{prop:t1-extremal}).
\end{enumerate}
\resultado \textbf{El EMV de la posición de Laplace es la mediana muestral} (cualquier punto entre los dos datos centrales si $n$ es par). $\blacksquare$
\comprobacion la densidad de Laplace tiene colas más pesadas que la normal, y su EMV es el estimador robusto frente a atípicos: el modelo elegido determina el estimador.
\end{solucion}

\begin{solucion}{eje:t4-regla3}
\planteamiento $\PP_p(\text{cero fallos})=(1-p)^{n}$, decreciente en $p$; usar $e^{x}\ge1+x$.
\pasos
\begin{enumerate}[label=\arabic*.]
  \item (a) $(1-p)^{n}\ge\alpha\iff1-p\ge\alpha^{1/n}\iff p\le1-\alpha^{1/n}$, pues $t\mapsto t^{n}$ es creciente en $[0,1]$. El conjunto es $[0,p_{\max}]$ con $p_{\max}=1-\alpha^{1/n}$.
  \item (b) $\alpha^{1/n}=e^{(\ln\alpha)/n}\ge1+\dfrac{\ln\alpha}{n}$, luego $p_{\max}=1-\alpha^{1/n}\le-\dfrac{\ln\alpha}{n}$. Para $\alpha=0{,}05$, $-\ln0{,}05\approx2{,}996\approx3$.
  \item (c) $n=300$: $p_{\max}=1-0{,}05^{1/300}\approx0{,}009\,936$; $3/n=0{,}01$.
\end{enumerate}
\resultado \textbf{$p_{\max}=1-\alpha^{1/n}\le-\ln(\alpha)/n\approx3/n$; con $n=300$, $p_{\max}\approx0{,}0099$}: con $95\,\%$ de confianza, la tasa de fallos graves es menor que un $1\,\%$. $\blacksquare$
\comprobacion $(1-0{,}009\,936)^{300}=e^{300\ln(0{,}990\,064)}\approx e^{-2{,}996}\approx0{,}05$.
\end{solucion}

\begin{solucion}{eje:t4-conservador}
\planteamiento calcular la cobertura aproximada con la normal y usar $\sqrt{p(1-p)}\le\frac12$.
\pasos
\begin{enumerate}[label=\arabic*.]
  \item Sea $c=z_{1-\alpha/2}$ y $\tau=\sqrt{p(1-p)/n}$. El intervalo cubre $p$ si y solo si $\abs{\widehat p-p}\le c/(2\sqrt n)$.
  \item Con $\widehat p\approx\Nor(p,\tau^{2})$: $\PP\bigl(\abs{\widehat p-p}\le\tfrac{c}{2\sqrt n}\bigr)\approx2\Phi\Bigl(\dfrac{c}{2\sqrt n\,\tau}\Bigr)-1=2\Phi\Bigl(\dfrac{c}{2\sqrt{p(1-p)}}\Bigr)-1$.
  \item Como $p(1-p)\le1/4$ con igualdad solo si $p=1/2$, el argumento es $\ge c$, y como $\Phi$ es creciente, la cobertura aproximada es $\ge2\Phi(c)-1=1-\alpha$, con igualdad solo si $p=1/2$.
  \item $p=0{,}1$: $\sqrt{p(1-p)}=0{,}3$ y el argumento vale $1{,}96/0{,}6\approx3{,}27$; cobertura aproximada $2\Phi(3{,}27)-1\approx0{,}9989$, independiente de $n$.
\end{enumerate}
\resultado \textbf{La cobertura aproximada es al menos $1-\alpha$}, y el intervalo es tanto más conservador cuanto más lejos está $p$ de $1/2$; con $p=0{,}1$ cubre el $99{,}9\,\%$ de las veces en lugar del $95\,\%$. $\blacksquare$
\comprobacion el precio es la amplitud: con $p=0{,}1$ el intervalo de Wald tendría un margen de $1{,}96\cdot0{,}3/\sqrt n$, un $40\,\%$ menor que el conservador $1{,}96\cdot0{,}5/\sqrt n$.
\end{solucion}

% ---------------------------------------------------------------------
\section{Tema 5}

\begin{solucion}{eje:t5-vf}
\planteamiento \cref{def:t5-pvalor}, \cref{prop:t5-pvalor} y \cref{obs:t5-pvalor}.
\pasos
\begin{enumerate}[label=(\alph*)]
  \item \textbf{Falsa.} El $p$-valor se calcula suponiendo $H_0$ cierta; no es $\PP(H_0\mid\text{datos})$.
  \item \textbf{Verdadera}: $0{,}03\le0{,}05$.
  \item \textbf{Falsa}: $0{,}03>0{,}01$.
  \item \textbf{Verdadera}: es la definición.
  \item \textbf{Falsa}: el $p$-valor no mide el tamaño del efecto; depende también de $n$.
  \item \textbf{Falsa}: no rechazar no demuestra $H_0$.
\end{enumerate}
\resultado \textbf{(a) F; (b) V; (c) F; (d) V; (e) F; (f) F.}
\end{solucion}

\begin{solucion}{eje:t5-moneda}
\datos $n=250$, $x=140$; bajo $H_0$, $X\sim\Bin(250;\allowbreak\,0{,}5)$, con media $125$ y desviación típica $\sqrt{62{,}5}\approx7{,}906$.
\planteamiento aproximación normal con corrección por continuidad (\cref{def:t3-cc}); contraste bilateral.
\pasos
\begin{enumerate}[label=\arabic*.]
  \item (a) $p=2\PP(X\ge140)\approx2\Bigl(1-\Phi\bigl(\tfrac{139{,}5-125}{7{,}906}\bigr)\Bigr)=2\bigl(1-\Phi(1{,}83)\bigr)=2\cdot0{,}0336=0{,}0672$.
  \item (b) $0{,}0672\le0{,}10$: se rechaza al $10\,\%$; $0{,}0672>0{,}05$: no se rechaza al $5\,\%$.
  \item (c) Al $1\,\%$ se rechaza si $\abs{x-125}-0{,}5\ge z_{0{,}995}\cdot7{,}906=2{,}576\cdot7{,}906\approx20{,}37$, es decir, $\abs{x-125}\ge20{,}87$: hacen falta al menos $146$ caras (o a lo sumo $104$).
  \item (d) Al $5\,\%$ se rechaza si $\abs{x-125}-0{,}5\ge1{,}96\cdot7{,}906\approx15{,}50$, es decir, $x\ge141$ o $x\le109$. Si $\theta=0{,}6$: media $150$, desviación típica $\sqrt{60}\approx7{,}746$, y $\PP(X\ge141)\approx1-\Phi\bigl(\tfrac{140{,}5-150}{7{,}746}\bigr)=\Phi(1{,}23)=0{,}8907$; $\PP(X\le109)$ es despreciable.
\end{enumerate}
\resultado \textbf{(a) $p\approx0{,}067$; (b) se rechaza al $10\,\%$, no al $5\,\%$; (c) $146$ caras; (d) potencia $\approx0{,}89$.}
\comprobacion los valores binomiales exactos son $p=0{,}0664$ (el «$7\,\%$» de la fuente), $2\PP(X\ge146)=0{,}0094\le0{,}01$ y potencia $0{,}8896$.
\end{solucion}

\begin{solucion}{eje:t5-poligrafo}
\datos sinceras: $131$ bien, $9$ mal; mentirosas: $15$ mal, $125$ bien.
\planteamiento con $H_0$: «dice la verdad», rechazar $H_0$ es clasificar como mentirosa.
\pasos
\begin{enumerate}[label=\arabic*.]
  \item Error de tipo~I: clasificar como mentirosa a una persona sincera. Estimación: $9/140\approx0{,}064$, que es el nivel estimado del detector.
  \item Error de tipo~II: clasificar como sincera a una persona que miente. Estimación: $15/140\approx0{,}107$.
  \item Potencia: $1-0{,}107=125/140\approx0{,}893$.
\end{enumerate}
\resultado \textbf{$\PP(\text{tipo I})\approx0{,}064$; $\PP(\text{tipo II})\approx0{,}107$; nivel $\approx0{,}064$; potencia $\approx0{,}893$.}
\comprobacion nivel $=\PP(\text{tipo I})$ y potencia $=1-\PP(\text{tipo II})$, como en la \cref{def:t5-errores}.
\end{solucion}

\begin{solucion}{eje:t5-radar}
\datos $n=3$; $\sigma=5$; $\alpha=0{,}04$; $\mu_0=40$.
\planteamiento contraste $z$ unilateral (\cref{prop:t5-potencia}); $\sigma/\sqrt3\approx2{,}887$; $z_{0{,}96}=1{,}7507$.
\pasos
\begin{enumerate}[label=\arabic*.]
  \item (a) $H_0:\mu=40$ (el coche no supera el límite; también puede tomarse la compuesta $\mu\le40$) frente a $H_1:\mu>40$. Se multa si $\xbar>40+1{,}7507\cdot2{,}887\approx45{,}05$.
  \item (b) Si nadie supera el límite, toda multa es un error de tipo~I: el $100\,\%$ de las multas serían erróneas (aunque solo se multaría a un $4\,\%$ o menos de los conductores).
  \item (c) $\pi(45)=1-\Phi\bigl(1{,}7507-5\sqrt3/5\bigr)=1-\Phi(1{,}7507-1{,}7321)=1-\Phi(0{,}019)\approx0{,}49$.
  \item (d) $n\ge\bigl((z_{0{,}96}+z_{0{,}90})\sigma/\delta\bigr)^{2}=\bigl((1{,}7507+1{,}2816)\cdot5/5\bigr)^{2}\approx9{,}19$, luego $10$ radares.
\end{enumerate}
\resultado \textbf{(a) umbral $\approx45{,}05$ km/h; (b) todas; (c) potencia $\approx0{,}49$; (d) $10$ radares.}
\comprobacion la potencia contra $45$ es casi $1/2$ porque el umbral ($45{,}05$) es casi igual a $45$: un coche a $45$ supera el umbral la mitad de las veces.
\end{solucion}

\begin{solucion}{eje:t5-uniforme}
\datos $X\sim\Unif(0,\theta)$; región de rechazo $\{x\le0{,}1\}\cup\{x\ge1{,}9\}$.
\planteamiento calcular probabilidades de intervalos con la densidad uniforme.
\pasos
\begin{enumerate}[label=\arabic*.]
  \item (a) Con $\theta=2$: $\PP(X\le0{,}1)+\PP(X\ge1{,}9)=0{,}1/2+0{,}1/2=0{,}1$.
  \item (b) Con $\theta=2{,}5$, no se rechaza si $0{,}1<X<1{,}9$: probabilidad $1{,}8/2{,}5=0{,}72$.
\end{enumerate}
\resultado \textbf{(a) $0{,}10$; (b) $0{,}72$.}
\comprobacion la potencia contra $\theta=2{,}5$ es $0{,}28$, mayor que el nivel $0{,}10$, pero pequeña: con una sola observación es difícil distinguir $\theta=2$ de $\theta=2{,}5$.
\end{solucion}

\begin{solucion}{eje:t5-t-test}
\datos $n=16$; $\xbar=11$; $s=2$; $\mu_0=10$; $\alpha=0{,}05$; $15$ grados de libertad.
\planteamiento contraste $t$; $t_{15;0{,}975}=2{,}131$ y $t_{15;0{,}95}=1{,}753$.
\pasos
\begin{enumerate}[label=\arabic*.]
  \item $t=\dfrac{11-10}{2/\sqrt{16}}=\dfrac{1}{0{,}5}=2$.
  \item (a) Bilateral: $\abs t=2<2{,}131$, no se rechaza. Como $1{,}753<2<2{,}131$, el $p$-valor unilateral está entre $0{,}025$ y $0{,}05$ y el bilateral entre $0{,}05$ y $0{,}10$.
  \item (b) Unilateral: $t=2>1{,}753$, se rechaza; $0{,}025<p<0{,}05$.
\end{enumerate}
\resultado \textbf{(a) No se rechaza ($0{,}05<p<0{,}10$); (b) se rechaza ($0{,}025<p<0{,}05$).}
\comprobacion los $p$-valores exactos son $0{,}0639$ y $0{,}0320$. El ejemplo muestra por qué la alternativa debe fijarse antes de ver los datos (\cref{obs:t5-errores}(a)).
\end{solucion}

\begin{solucion}{eje:t5-openstax-z}
\datos $\mu_0=16{,}43$; $\sigma=0{,}8$; $n=15$; $\xbar=16$; $\alpha=0{,}05$.
\planteamiento contraste $z$ unilateral: $H_0:\mu=16{,}43$ frente a $H_1:\mu<16{,}43$.
\pasos
\begin{enumerate}[label=\arabic*.]
  \item $\sigma/\sqrt n=0{,}8/\sqrt{15}\approx0{,}2066$.
  \item $z=(16-16{,}43)/0{,}2066\approx-2{,}08$.
  \item $p=\Phi(-2{,}08)=1-0{,}9812=0{,}0188<0{,}05$.
\end{enumerate}
\resultado \textbf{$z\approx-2{,}08$, $p\approx0{,}019$: se rechaza $H_0$}; hay evidencia de que la media es menor que la anunciada.
\comprobacion equivalentemente, $z=-2{,}08<-z_{0{,}95}=-1{,}645$.
\end{solucion}

\begin{solucion}{eje:t5-openstax-prop}
\datos $p_0=0{,}25$; $n=42$; $x=17$; $\alpha=0{,}01$.
\planteamiento contraste $z$ unilateral para una proporción, con el error estándar calculado con $p_0$.
\pasos
\begin{enumerate}[label=\arabic*.]
  \item $\widehat p=17/42\approx0{,}4048$.
  \item $\sqrt{p_0(1-p_0)/n}=\sqrt{0{,}1875/42}\approx0{,}06682$.
  \item $z=(0{,}4048-0{,}25)/0{,}06682\approx2{,}32$; $p=1-\Phi(2{,}32)=1-0{,}9898=0{,}0102$.
  \item $0{,}0102>0{,}01$: no se rechaza.
\end{enumerate}
\resultado \textbf{$z\approx2{,}32$, $p\approx0{,}010$: no se rechaza $H_0$ al $1\,\%$} (sí se rechazaría al $5\,\%$).
\comprobacion $np_0=10{,}5$ y $n(1-p_0)=31{,}5$, de modo que la aproximación normal es aceptable; el resultado está justo en el límite, y conviene informar del $p$-valor en lugar de solo «no significativo».
\end{solucion}

\begin{solucion}{eje:t5-chi-dias}
\datos $O=(15,12,9,9,15)$; $n=60$; $k=5$.
\planteamiento \cref{def:t5-chi2} con $p_{0j}=1/5$.
\pasos
\begin{enumerate}[label=\arabic*.]
  \item $E_j=60/5=12$ para todo $j$.
  \item $X^{2}=\dfrac{3^{2}+0^{2}+3^{2}+3^{2}+3^{2}}{12}=\dfrac{36}{12}=3$.
  \item Grados de libertad: $k-1=4$; $\chi^{2}_{4;0{,}95}=9{,}488$.
  \item $3<9{,}488$: no se rechaza.
\end{enumerate}
\resultado \textbf{$X^{2}=3$ con $4$ grados de libertad; no se rechaza la uniformidad ($p\approx0{,}56$).}
\comprobacion $\E[\chi^{2}_4]=4$: un valor de $3$ es completamente típico bajo $H_0$.
\end{solucion}

\begin{solucion}{eje:t5-benford}
\datos $n=100$; $E_j=100p_j$; $O=(7,13,12,9,9,13,11,10,16)$.
\planteamiento \cref{def:t5-chi2} con $k-1=8$ grados de libertad.
\pasos
\begin{enumerate}[label=\arabic*.]
  \item $E=(30{,}1;\allowbreak\,17{,}6;\allowbreak\,12{,}5;\allowbreak\,9{,}7;\allowbreak\,7{,}9;\allowbreak\,6{,}7;\allowbreak\,5{,}8;\allowbreak\,5{,}1;\allowbreak\,4{,}6)$.
  \item Contribuciones $(O_j-E_j)^{2}/E_j$: $17{,}73$; $1{,}20$; $0{,}02$; $0{,}05$; $0{,}15$; $5{,}92$; $4{,}66$; $4{,}71$; $28{,}25$.
  \item $X^{2}\approx62{,}70$; $\chi^{2}_{8;0{,}999}=26{,}12$.
  \item $62{,}70>26{,}12$: se rechaza $H_0$ al nivel $0{,}001$.
\end{enumerate}
\resultado \textbf{$X^{2}\approx62{,}7$ con $8$ grados de libertad: se rechaza con mucha claridad} (hay demasiados nueves y muy pocos unos); se recomienda revisar la nota de gastos.
\comprobacion $E_9=4{,}6<5$. Agrupando los dígitos $8$ y $9$ ($E=9{,}7$, $O=26$) se obtiene $X^{2}\approx57{,}1$ con $7$ grados de libertad, frente a $\chi^{2}_{7;0{,}999}=24{,}32$: la conclusión no cambia.
\end{solucion}

\begin{solucion}{eje:t5-dos-modelos}
\datos $n_A=n_B=1000$; $\widehat p_A=0{,}78$; $\widehat p_B=0{,}75$.
\planteamiento bajo $H_0$, $p_A=p_B=p$, $\Var(\widehat p_A-\widehat p_B)=p(1-p)(1/n_A+1/n_B)$, y $p$ se estima con la proporción combinada.
\pasos
\begin{enumerate}[label=\arabic*.]
  \item (a) $H_0:p_A=p_B$ (compuesta: no fija el valor común) frente a $H_1:p_A\ne p_B$.
  \item (b) $\widehat p=(780+750)/2000=0{,}765$; $\widehat{\Var}=0{,}765\cdot0{,}235\cdot(2/1000)\approx0{,}000\,359\,6$; $\widehat{\ee}\approx0{,}01896$.
  \item $z=(0{,}78-0{,}75)/0{,}01896\approx1{,}58$; $p=2\bigl(1-\Phi(1{,}58)\bigr)=2\cdot0{,}0571=0{,}1142$.
  \item (c) $0{,}114>0{,}05$: no es significativa.
\end{enumerate}
\resultado \textbf{$z\approx1{,}58$, $p\approx0{,}11$: la diferencia de tres puntos no es significativa al $5\,\%$.}
\comprobacion para que una diferencia de tres puntos fuera significativa con estas exactitudes haría falta que $0{,}03/\sqrt{0{,}18\cdot2/n}\ge1{,}96$, es decir, $n\ge1{,}96^{2}\cdot0{,}36/0{,}0009\approx1537$ preguntas por modelo.
\end{solucion}

\begin{solucion}{eje:t5-multiples}
\datos $m=20$; $\alpha=0{,}05$; todas las nulas ciertas; contrastes independientes.
\planteamiento \cref{prop:t5-multiples} y \cref{teo:t5-bonferroni}.
\pasos
\begin{enumerate}[label=\arabic*.]
  \item (a) $\PP(\text{al menos uno})=1-0{,}95^{20}\approx1-0{,}3585=0{,}6415$. Número esperado: $20\cdot0{,}05=1$.
  \item (b) Bonferroni: nivel $0{,}05/20=0{,}0025$ por contraste.
  \item Probabilidad exacta de algún falso positivo con contrastes independientes: $1-(1-0{,}0025)^{20}=1-0{,}9975^{20}\approx0{,}0488$.
\end{enumerate}
\resultado \textbf{(a) $\approx0{,}64$ y $1$ falso positivo esperado; (b) $0{,}0025$ por contraste y FWER $\approx0{,}049\le0{,}05$.}
\comprobacion el valor $0{,}0488$ es algo menor que $0{,}05$, como garantiza el \cref{teo:t5-bonferroni}; la diferencia es la holgura de la desigualdad de la unión.
\end{solucion}

\begin{solucion}{eje:t5-sidak}
\planteamiento (a) complemento e independencia; (b) igualar a $\alpha$; (c) desigualdad de Bernoulli.
\pasos
\begin{enumerate}[label=\arabic*.]
  \item (a) $\PP(\min_iP_i>t)=\PP(P_1>t,\dots,P_n>t)=\prod_i\PP(P_i>t)=(1-t)^{n}$, luego $\PP(\min_iP_i\le t)=1-(1-t)^{n}$.
  \item (b) Se produce algún falso positivo si y solo si $\min_iP_i\le t$. Con $t=1-(1-\alpha)^{1/n}$: $1-(1-t)^{n}=1-\bigl((1-\alpha)^{1/n}\bigr)^{n}=1-(1-\alpha)=\alpha$.
  \item (c) Con $x=-\alpha/n\ge-1$, Bernoulli da $(1-\alpha/n)^{n}\ge1-\alpha$. Tomando raíces $n$-ésimas (función creciente): $1-\alpha/n\ge(1-\alpha)^{1/n}$, es decir, $\alpha/n\le1-(1-\alpha)^{1/n}$.
  \item Un umbral mayor rechaza más hipótesis, luego el procedimiento es al menos tan potente como el de Bonferroni.
\end{enumerate}
\resultado \textbf{$\PP(\min P_i\le t)=1-(1-t)^{n}$; el umbral $1-(1-\alpha)^{1/n}$ da FWER exactamente $\alpha$ y es mayor o igual que $\alpha/n$.} $\blacksquare$
\comprobacion con $n=20$ y $\alpha=0{,}05$: $1-0{,}95^{1/20}\approx0{,}002\,561>0{,}0025$. La ganancia es pequeña, y a cambio exige independencia; Bonferroni no la exige.
\end{solucion}

\begin{solucion}{eje:t5-uniforme-p}
\planteamiento calcular la función de distribución de $p$ usando que $F_0$ es continua y estrictamente creciente, y por tanto invertible en su recorrido.
\pasos
\begin{enumerate}[label=\arabic*.]
  \item Sea $u\in(0,1)$. $p\le u\iff1-F_0(T)\le u\iff F_0(T)\ge1-u\iff T\ge F_0^{-1}(1-u)$.
  \item Bajo $H_0$, $\PP\bigl(T\ge F_0^{-1}(1-u)\bigr)=1-F_0\bigl(F_0^{-1}(1-u)\bigr)=1-(1-u)=u$ (por continuidad, $\PP(T=c)=0$).
  \item Luego $\PP_{H_0}(p\le u)=u$ para todo $u\in(0,1)$: $p\sim\Unif(0,1)$. Con $u=\alpha$, $\PP_{H_0}(p\le\alpha)=\alpha$.
\end{enumerate}
\resultado \textbf{Bajo $H_0$, el $p$-valor es uniforme en $(0,1)$ y la regla «rechazar si $p\le\alpha$» tiene nivel exactamente $\alpha$.} $\blacksquare$
\comprobacion este resultado justifica el apartado (a) del \cref{eje:t5-sidak}. Si el estadístico es discreto (por ejemplo, binomial), $p$ no es exactamente uniforme y el nivel real es $\le\alpha$.
\end{solucion}

\begin{solucion}{eje:t5-dualidad}
\planteamiento escribir la condición de no rechazo y despejar $\mu_0$.
\pasos
\begin{enumerate}[label=\arabic*.]
  \item Sea $c=z_{1-\alpha/2}$. El contraste no rechaza $H_0:\mu=\mu_0$ si y solo si $\abs{\xbar-\mu_0}\le c\,\sigma/\sqrt n$.
  \item $\abs{\xbar-\mu_0}\le c\sigma/\sqrt n\iff-c\sigma/\sqrt n\le\mu_0-\xbar\le c\sigma/\sqrt n\iff\xbar-c\sigma/\sqrt n\le\mu_0\le\xbar+c\sigma/\sqrt n$.
  \item La última condición es la pertenencia de $\mu_0$ al intervalo del \cref{teo:t4-ic-z}.
\end{enumerate}
\resultado \textbf{No se rechaza $\mu_0$ al nivel $\alpha$ si y solo si $\mu_0$ está en el intervalo de nivel $1-\alpha$.} El intervalo de confianza es el conjunto de valores del parámetro compatibles con los datos al nivel $\alpha$. $\blacksquare$
\comprobacion en el \cref{eje:t4-interpretacion}, el $0$ no está en $[10,20]$, luego un contraste bilateral al $5\,\%$ rechaza la mejora nula.
\end{solucion}

\begin{solucion}{eje:t5-potencia}
\datos $\pi(\mu)=1-\Phi\bigl(z_{1-\alpha}-(\mu-\mu_0)\sqrt n/\sigma\bigr)$.
\planteamiento monotonía de $\Phi$; resolver la desigualdad $\pi(\mu_1)\ge1-\beta$.
\pasos
\begin{enumerate}[label=\arabic*.]
  \item (a) El argumento $g(\mu)=z_{1-\alpha}-(\mu-\mu_0)\sqrt n/\sigma$ es estrictamente decreciente en $\mu$ y $\Phi$ es estrictamente creciente, luego $\pi=1-\Phi\circ g$ es estrictamente creciente. Además $\pi(\mu_0)=1-\Phi(z_{1-\alpha})=1-(1-\alpha)=\alpha$.
  \item (b) Si $\mu>\mu_0$, $g(\mu)\to-\infty$ cuando $n\to\infty$, y $\Phi(g)\to0$, luego $\pi(\mu)\to1$.
  \item (c) $\pi(\mu_1)\ge1-\beta\iff\Phi\bigl(z_{1-\alpha}-\delta\sqrt n/\sigma\bigr)\le\beta\iff z_{1-\alpha}-\delta\sqrt n/\sigma\le z_\beta=-z_{1-\beta}\iff\sqrt n\ge(z_{1-\alpha}+z_{1-\beta})\sigma/\delta$, y se elevan al cuadrado ambos miembros, positivos.
  \item (d) $n\ge\bigl((1{,}645+1{,}2816)\cdot80/20\bigr)^{2}=(11{,}706)^{2}\approx137{,}0$, luego $n=138$.
\end{enumerate}
\resultado \textbf{$\pi$ es creciente, vale $\alpha$ en $\mu_0$ y tiende a $1$; el tamaño muestral es el indicado; en el ejemplo, $n=138$ respuestas.} $\blacksquare$
\comprobacion con $n=138$: $\pi(220)=1-\Phi\bigl(1{,}645-20\sqrt{138}/80\bigr)=1-\Phi(1{,}645-2{,}937)=\Phi(1{,}292)\approx0{,}90$.
\end{solucion}

% ---------------------------------------------------------------------
\section{Tema 6}

\begin{solucion}{eje:t6-vf}
\planteamiento \cref{teo:t6-mc}, \cref{prop:t6-residuos}, \cref{cor:t6-r2} y \cref{obs:t6-errores}.
\pasos
\begin{enumerate}[label=(\alph*)]
  \item \textbf{Verdadera}: $\widehat b_0=\ybar-\widehat b_1\xbar$ equivale a $\ybar=\widehat b_0+\widehat b_1\xbar$.
  \item \textbf{Falsa}: $R^{2}=r^{2}=0{,}81\ge0$.
  \item \textbf{Falsa}: $R^{2}$ mide la calidad del ajuste lineal, no la causalidad.
  \item \textbf{Verdadera}, por la \cref{prop:t6-residuos}.
  \item \textbf{Falsa}: la pendiente de $x$ sobre $y$ es $r\,s_x/s_y$, que coincide con $1/\widehat b_1=s_x/(r\,s_y)$ solo si $r^{2}=1$.
\end{enumerate}
\resultado \textbf{(a) V; (b) F; (c) F; (d) V; (e) F.}
\end{solucion}

\begin{solucion}{eje:t6-basico}
\datos $n=5$; $\sum x_i=10$; $\sum y_i=15$; $\sum x_i^{2}=30$; $\sum y_i^{2}=55$; $\sum x_iy_i=38$.
\planteamiento sumas centradas con la fórmula de cálculo y \cref{teo:t6-mc}.
\pasos
\begin{enumerate}[label=\arabic*.]
  \item $\xbar=2$, $\ybar=3$; $S_{xx}=30-5\cdot4=10$; $S_{yy}=55-5\cdot9=10$; $S_{xy}=38-5\cdot2\cdot3=8$; $r=8/\sqrt{100}=0{,}8$.
  \item $\widehat b_1=8/10=0{,}8$; $\widehat b_0=3-0{,}8\cdot2=1{,}4$; $\widehat y=1{,}4+0{,}8x$.
  \item Ajustados: $1{,}4;\allowbreak\ 2{,}2;\allowbreak\ 3{,}0;\allowbreak\ 3{,}8;\allowbreak\ 4{,}6$. Residuos: $-0{,}4;\allowbreak\ 0{,}8;\allowbreak\ -1{,}0;\allowbreak\ 1{,}2;\allowbreak\ -0{,}6$.
  \item $\operatorname{SC}_E=0{,}16+0{,}64+1+1{,}44+0{,}36=3{,}6$; $R^{2}=1-3{,}6/10=0{,}64$; $s_e=\sqrt{3{,}6/3}=\sqrt{1{,}2}\approx1{,}095$.
  \item $\widehat y(2{,}5)=1{,}4+2=3{,}4$.
\end{enumerate}
\resultado \textbf{$\widehat y=1{,}4+0{,}8x$; $r=0{,}8$; $R^{2}=0{,}64$; $s_e\approx1{,}10$; predicción $3{,}4$.}
\comprobacion los residuos suman $0$ y $\operatorname{SC}_E=S_{yy}(1-r^{2})=10\cdot0{,}36=3{,}6$.
\end{solucion}

\begin{solucion}{eje:t6-gauss}
\datos $(1,8),(3,2),(5,1)$.
\planteamiento la log-verosimilitud total es una constante menos $\operatorname{SC}(b,a)/(2\sigma^{2})$; se maximiza minimizando $\operatorname{SC}$ (\cref{obs:t6-modelo}).
\pasos
\begin{enumerate}[label=\arabic*.]
  \item (a) $y_1\sim\Nor(ax_1+b,\sigma^{2})$: $\ln f(y_1)=-\ln\bigl(\sigma\sqrt{2\pi}\bigr)-\dfrac{(y_1-ax_1-b)^{2}}{2\sigma^{2}}$.
  \item (b) Sumando para los tres datos, $\ell(a,b)=-3\ln(\sigma\sqrt{2\pi})-\frac{1}{2\sigma^{2}}\sum_i(y_i-ax_i-b)^{2}$; para cualquier $\sigma$, el máximo en $(a,b)$ es el mínimo de la suma de cuadrados.
  \item $\xbar=3$; $\ybar=11/3$; $S_{xx}=35-27=8$; $S_{xy}=19-3\cdot3\cdot\frac{11}3=19-33=-14$.
  \item $\widehat a=-14/8=-1{,}75$; $\widehat b=\frac{11}{3}+1{,}75\cdot3=\frac{107}{12}\approx8{,}917$.
\end{enumerate}
\resultado \textbf{$\widehat a=-1{,}75$ y $\widehat b=107/12\approx8{,}92$}, independientes de $\sigma$.
\comprobacion residuos: $8-7{,}167=0{,}833$; $2-3{,}667=-1{,}667$; $1-0{,}167=0{,}833$; suman $0$ y $\sum x_ie_i=0{,}833-5+4{,}167=0$.
\end{solucion}

\begin{solucion}{eje:t6-examenes}
\datos $n=11$ y las sumas del enunciado.
\planteamiento fórmula de cálculo para las sumas centradas y \cref{teo:t6-mc}.
\pasos
\begin{enumerate}[label=\arabic*.]
  \item $\xbar=761/11\approx69{,}18$; $\ybar=1765/11\approx160{,}45$.
  \item $S_{xx}=52\,729-761^{2}/11\approx52\,729-52\,647{,}36=81{,}64$.
  \item $S_{xy}=122\,500-761\cdot1765/11\approx122\,500-122\,105{,}91=394{,}09$.
  \item $S_{yy}=287\,529-1765^{2}/11\approx287\,529-283\,202{,}27=4326{,}73$.
  \item $\widehat b_1=394{,}09/81{,}64\approx4{,}827$; $\widehat b_0=160{,}45-4{,}827\cdot69{,}18\approx-173{,}51$.
  \item $r=394{,}09/\sqrt{81{,}64\cdot4326{,}73}\approx394{,}09/594{,}31\approx0{,}663$.
  \item (b) $\widehat y(73)=-173{,}51+4{,}827\cdot73\approx178{,}9$.
  \item (c) $\widehat y(65)\approx-173{,}51+313{,}78=140{,}27$; residuo $175-140{,}27\approx34{,}7$.
\end{enumerate}
\resultado \textbf{$\widehat y=-173{,}51+4{,}83x$; $r\approx0{,}663$; predicción $\approx178{,}9$; residuo $\approx34{,}7$.}
\comprobacion coincide con la recta y la correlación de la fuente; $R^{2}\approx0{,}44$: el tercer examen explica menos de la mitad de la variabilidad del final.
\end{solucion}

\begin{solucion}{eje:t6-buceo}
\datos $n=6$; $\sum x_i=450$; $\sum y_i=262$; $\sum x_i^{2}=35\,500$; $\sum y_i^{2}=13\,784$; $\sum x_iy_i=17\,700$.
\planteamiento \cref{teo:t6-mc}; después, examinar los residuos.
\pasos
\begin{enumerate}[label=\arabic*.]
  \item $\xbar=75$; $\ybar\approx43{,}67$; $S_{xx}=35\,500-33\,750=1750$; $S_{xy}=17\,700-450\cdot262/6=17\,700-19\,650=-1950$; $S_{yy}=13\,784-11\,440{,}67=2343{,}33$.
  \item $\widehat b_1=-1950/1750\approx-1{,}114$; $\widehat b_0=43{,}67+1{,}114\cdot75\approx127{,}24$.
  \item $r=-1950/\sqrt{1750\cdot2343{,}33}\approx-1950/2025{,}0\approx-0{,}963$.
  \item (b) $\widehat y(85)\approx127{,}24-94{,}71=32{,}5$ minutos; $\widehat y(110)\approx127{,}24-122{,}57=4{,}7$ minutos.
  \item La primera es una interpolación (dentro de $[50,100]$) y es razonable; la segunda es una extrapolación y no es fiable.
  \item (c) Residuos: $8{,}5;\allowbreak\ -5{,}4;\allowbreak\ -4{,}2;\allowbreak\ -3{,}1;\allowbreak\ -2{,}0;\allowbreak\ 6{,}2$. Son positivos en los extremos y negativos en el centro: la relación es curva (convexa), y la recta subestima en los extremos.
\end{enumerate}
\resultado \textbf{$\widehat y=127{,}24-1{,}114x$; $r\approx-0{,}963$; predicciones $32{,}5$ (fiable) y $4{,}7$ (extrapolación); los residuos sugieren curvatura.}
\comprobacion los residuos suman $8{,}48-5{,}38-4{,}24-3{,}10-1{,}95+6{,}19\approx0$. A pesar de $R^{2}\approx0{,}93$, el patrón de residuos indica que un modelo curvo sería mejor (\cref{obs:t6-errores}(e)).
\end{solucion}

\begin{solucion}{eje:t6-helados}
\datos $n=9$ y las sumas del enunciado.
\planteamiento \cref{teo:t6-mc} y \cref{cor:t6-r2}.
\pasos
\begin{enumerate}[label=\arabic*.]
  \item $\xbar=654/9\approx72{,}67$; $\ybar=1288/9\approx143{,}11$.
  \item $S_{xx}=47\,988-654^{2}/9=47\,988-47\,524=464$.
  \item $S_{xy}=98\,132-654\cdot1288/9\approx98\,132-93\,594{,}67=4537{,}33$.
  \item $S_{yy}=231\,026-1288^{2}/9\approx231\,026-184\,327{,}11=46\,698{,}89$.
  \item $\widehat b_1=4537{,}33/464\approx9{,}779$; $\widehat b_0=143{,}11-9{,}779\cdot72{,}67\approx-567{,}48$.
  \item $r=4537{,}33/\sqrt{464\cdot46\,698{,}89}\approx4537{,}33/4654{,}9\approx0{,}975$; $R^{2}\approx0{,}950$.
  \item $\widehat y(74)\approx-567{,}48+9{,}779\cdot74\approx156{,}1$.
\end{enumerate}
\resultado \textbf{$\widehat y=-567{,}48+9{,}78x$; $r\approx0{,}975$; $R^{2}\approx0{,}95$; unos $156$ helados.}
\comprobacion la ordenada en el origen, $-567$, no tiene interpretación (no se venden helados negativos a $0$ grados): $x=0$ está muy fuera del rango de los datos.
\end{solucion}

\begin{solucion}{eje:t6-resumen}
\datos $\xbar=10$; $\ybar=50$; $s_x=2$; $s_y=8$; $r=0{,}75$.
\planteamiento \cref{teo:t6-mc} con $\widehat b_1=r\,s_y/s_x$; para $x$ sobre $y$, intercambiar los papeles.
\pasos
\begin{enumerate}[label=\arabic*.]
  \item (a) $\widehat b_1=0{,}75\cdot8/2=3$; $\widehat b_0=50-3\cdot10=20$; $\widehat y=20+3x$; $\widehat y(12)=56$.
  \item (b) Pendiente $r\,s_x/s_y=0{,}75\cdot2/8=0{,}1875$; ordenada $10-0{,}1875\cdot50=0{,}625$; $\widehat x=0{,}625+0{,}1875\,y$.
  \item (c) $3\cdot0{,}1875=0{,}5625=0{,}75^{2}$. Despejando $x$ en la recta de (a) se obtendría $x=(y-20)/3$, de pendiente $1/3\neq0{,}1875$.
\end{enumerate}
\resultado \textbf{$\widehat y=20+3x$ (predicción $56$) y $\widehat x=0{,}625+0{,}1875y$; el producto de pendientes es $r^{2}=0{,}5625$.}
\comprobacion ambas rectas pasan por $(10,50)$: $20+30=50$ y $0{,}625+9{,}375=10$.
\end{solucion}

\begin{solucion}{eje:t6-loglog}
\datos $u_i=\log_{10}N_i=6,7,8,9$; $v_i=\log_{10}L_i=0{,}6990;\allowbreak\ 0{,}6128;\allowbreak\ 0{,}5185;\allowbreak\ 0{,}4314$.
\planteamiento recta de mínimos cuadrados de $v$ sobre $u$; $\widehat\alpha=-\widehat b_1$ y $\widehat a=10^{\widehat b_0}$.
\pasos
\begin{enumerate}[label=\arabic*.]
  \item $\bar u=7{,}5$; $\bar v=2{,}2617/4=0{,}565\,43$; $S_{uu}=2{,}25+0{,}25+0{,}25+2{,}25=5$.
  \item $S_{uv}=\sum(u_i-\bar u)v_i=-1{,}5\cdot0{,}6990-0{,}5\cdot0{,}6128+0{,}5\cdot0{,}5185+1{,}5\cdot0{,}4314=-0{,}448\,55$.
  \item $\widehat b_1=-0{,}448\,55/5\approx-0{,}0897$; $\widehat b_0=0{,}565\,43+0{,}0897\cdot7{,}5\approx1{,}2382$.
  \item $\widehat\alpha\approx0{,}0897$; $\widehat a=10^{1{,}2382}\approx17{,}3$; $L\approx17{,}3\,N^{-0{,}0897}$.
  \item (b) $\log_{10}\widehat L=1{,}2382-0{,}0897\cdot10\approx0{,}3411$; $\widehat L\approx10^{0{,}3411}\approx2{,}19$.
  \item (c) $N=10^{10}$ está fuera del rango observado ($10^{6}$ a $10^{9}$); la predicción supone que la ley de potencia sigue valiendo un orden de magnitud más allá.
\end{enumerate}
\resultado \textbf{$L\approx17{,}3\,N^{-0{,}090}$; predicción para $N=10^{10}$: $L\approx2{,}19$ (extrapolación).}
\comprobacion con la ley ajustada, $L(10^{6})=17{,}3\cdot10^{-0{,}538}\approx17{,}3\cdot0{,}290\approx5{,}02$, muy próximo al dato $5{,}0$; $r\approx-0{,}9999$.
\end{solucion}

\begin{solucion}{eje:t6-simpson}
\datos A: $45/50$ fáciles, $75/250$ difíciles. B: $200/250$ fáciles, $10/50$ difíciles.
\planteamiento calcular proporciones por grupo y en total, y observar la composición de cada evaluación.
\pasos
\begin{enumerate}[label=\arabic*.]
  \item A: fáciles $0{,}90$; difíciles $0{,}30$; total $120/300=0{,}40$.
  \item B: fáciles $0{,}80$; difíciles $0{,}20$; total $210/300=0{,}70$.
  \item A es mejor en cada tipo de pregunta, pero peor en el total: A se ha evaluado sobre todo con preguntas difíciles ($250$ de $300$) y B sobre todo con fáciles. La dificultad es una variable de confusión.
  \item Una comparación justa usa la misma mezcla para ambos. Con una mezcla al $50\,\%$: A obtiene $(0{,}90+0{,}30)/2=0{,}60$ y B, $(0{,}80+0{,}20)/2=0{,}50$.
\end{enumerate}
\resultado \textbf{A es mejor} (lo es en ambos estratos); la exactitud total favorece a B solo porque los dos modelos se evaluaron con mezclas distintas de dificultad (paradoja de Simpson).
\comprobacion con cualquier mezcla común con proporción $w$ de fáciles, A obtiene $0{,}30+0{,}60w$ y B, $0{,}20+0{,}60w$: A gana por $0{,}10$ siempre.
\end{solucion}

\begin{solucion}{eje:t6-cuadratica}
\datos $x=-2,\dots,2$; $y=x^{2}$.
\planteamiento calcular $S_{xy}$ y los residuos.
\pasos
\begin{enumerate}[label=\arabic*.]
  \item $\xbar=0$; $\ybar=10/5=2$; $S_{xy}=\sum x_iy_i-0=\sum x_i^{3}=-8-1+0+1+8=0$. Luego $r=0$.
  \item $\widehat b_1=0$; $\widehat b_0=2$; la recta es $\widehat y=2$.
  \item Residuos: $2,-1,-2,-1,2$. Forman una «U»: positivos en los extremos, negativos en el centro.
\end{enumerate}
\resultado \textbf{$r=0$ y la recta es horizontal, pero $y$ está determinada exactamente por $x$}; el patrón de los residuos revela la relación cuadrática.
\comprobacion $\operatorname{SC}_E=4+1+4+1+4=14=S_{yy}$, coherente con $R^{2}=0$.
\end{solucion}

\begin{solucion}{eje:t6-normales}
\planteamiento derivadas parciales de $\operatorname{SC}(b_0,b_1)=\sum_i(y_i-b_0-b_1x_i)^{2}$.
\pasos
\begin{enumerate}[label=\arabic*.]
  \item $\dfrac{\partial\operatorname{SC}}{\partial b_0}=-2\sum_i(y_i-b_0-b_1x_i)$ y $\dfrac{\partial\operatorname{SC}}{\partial b_1}=-2\sum_ix_i(y_i-b_0-b_1x_i)$.
  \item Igualando a cero se obtienen las ecuaciones normales: $nb_0+b_1\sum x_i=\sum y_i$ y $b_0\sum x_i+b_1\sum x_i^{2}=\sum x_iy_i$.
  \item De la primera, $b_0=\ybar-b_1\xbar$. Sustituyendo en la segunda: $(\ybar-b_1\xbar)n\xbar+b_1\sum x_i^{2}=\sum x_iy_i$, es decir, $b_1\bigl(\sum x_i^{2}-n\xbar^{2}\bigr)=\sum x_iy_i-n\xbar\ybar$, y $b_1=S_{xy}/S_{xx}$.
  \item Hessiana: $H=\begin{pmatrix}2n&2\sum x_i\\2\sum x_i&2\sum x_i^{2}\end{pmatrix}$, constante. Su elemento $(1,1)$ es $2n>0$ y su determinante es $4\bigl(n\sum x_i^{2}-(\sum x_i)^{2}\bigr)=4nS_{xx}>0$. Por el criterio de Sylvester, $H$ es definida positiva.
  \item Una función cuadrática con hessiana constante definida positiva es estrictamente convexa, y su único punto crítico es el mínimo global.
\end{enumerate}
\resultado \textbf{Se recupera $\widehat b_1=S_{xy}/S_{xx}$ y $\widehat b_0=\ybar-\widehat b_1\xbar$, y el punto crítico es el mínimo global.} $\blacksquare$
\comprobacion con los datos del \cref{ej:t6-recta}: $5b_0+15b_1=20$ y $15b_0+55b_1=66$ tienen solución $b_1=0{,}6$, $b_0=2{,}2$.
\end{solucion}

\begin{solucion}{eje:t6-residuos}
\planteamiento las ecuaciones normales del \cref{eje:t6-normales} se escriben en términos de los residuos $e_i=y_i-\widehat b_0-\widehat b_1x_i$.
\pasos
\begin{enumerate}[label=\arabic*.]
  \item La primera ecuación normal es $\sum_i(y_i-\widehat b_0-\widehat b_1x_i)=0$, es decir, $\sum_ie_i=0$. (Directamente: $e_i=(y_i-\ybar)-\widehat b_1(x_i-\xbar)$ y ambas desviaciones suman cero.)
  \item La segunda es $\sum_ix_i(y_i-\widehat b_0-\widehat b_1x_i)=0$, es decir, $\sum_ix_ie_i=0$. (Directamente: $\sum_ix_ie_i=\sum_i(x_i-\xbar)e_i=S_{xy}-\widehat b_1S_{xx}=0$.)
  \item $\sum_i\widehat y_ie_i=\widehat b_0\sum_ie_i+\widehat b_1\sum_ix_ie_i=0$.
  \item Como $y_i=\widehat y_i+e_i$ y $\sum e_i=0$, $\frac1n\sum\widehat y_i=\ybar$.
\end{enumerate}
\resultado \textbf{Las tres identidades son ciertas, y la media de los ajustados es $\ybar$.} $\blacksquare$
\comprobacion en el \cref{ej:t6-recta}: $\sum x_ie_i=-0{,}8+1{,}2+3{,}0-2{,}4-1{,}0=0$.
\end{solucion}

\begin{solucion}{eje:t6-origen}
\planteamiento derivar $g(b)=\sum_i(y_i-bx_i)^{2}$.
\pasos
\begin{enumerate}[label=\arabic*.]
  \item $g'(b)=-2\sum_ix_i(y_i-bx_i)=0\iff\sum x_iy_i=b\sum x_i^{2}\iff b=\sum x_iy_i/\sum x_i^{2}$.
  \item $g''(b)=2\sum x_i^{2}>0$: $g$ es una parábola convexa y el punto crítico es su mínimo global.
  \item La ecuación $g'(\widehat b)=0$ es precisamente $\sum_ix_ie_i=0$.
  \item Ejemplo: $(1,2),(2,2)$. $\widehat b=(2+4)/(1+4)=1{,}2$; residuos $2-1{,}2=0{,}8$ y $2-2{,}4=-0{,}4$; $\sum e_i=0{,}4\ne0$, mientras que $\sum x_ie_i=0{,}8-0{,}8=0$.
\end{enumerate}
\resultado \textbf{$\widehat b=\sum x_iy_i/\sum x_i^{2}$; $\sum x_ie_i=0$, pero en general $\sum e_i\neq0$}, porque sin ordenada en el origen no hay ecuación normal asociada a $b_0$. $\blacksquare$
\comprobacion con ordenada libre, esos dos puntos dan la recta horizontal $y=2$ con residuos nulos.
\end{solucion}

\begin{solucion}{eje:t6-transformaciones}
\planteamiento por la \cref{prop:t1-lineal}, $x_i'-\xbar'=\alpha(x_i-\xbar)$ e $y_i'-\ybar'=\gamma(y_i-\ybar)$.
\pasos
\begin{enumerate}[label=\arabic*.]
  \item $S_{x'y'}=\alpha\gamma S_{xy}$ y $S_{x'x'}=\alpha^{2}S_{xx}$, luego $\widehat b_1'=\alpha\gamma S_{xy}/(\alpha^{2}S_{xx})=(\gamma/\alpha)\widehat b_1$.
  \item Por el \cref{eje:t1-r-invariante}, $r'=\operatorname{signo}(\alpha\gamma)\,r$, luego $R^{2}=r'^{2}=r^{2}$.
  \item Recta ajustada: $\widehat y_i'-\ybar'=\widehat b_1'(x_i'-\xbar')=(\gamma/\alpha)\widehat b_1\,\alpha(x_i-\xbar)=\gamma(\widehat y_i-\ybar)$, de modo que $\widehat y_i'=\gamma\widehat y_i+\delta$.
  \item Residuos: $e_i'=y_i'-\widehat y_i'=\gamma y_i+\delta-\gamma\widehat y_i-\delta=\gamma e_i$.
  \item De metros a centímetros: $\alpha=100$, $\gamma=1$; la pendiente se divide por $100$.
\end{enumerate}
\resultado \textbf{$\widehat b_1'=(\gamma/\alpha)\widehat b_1$, $R^{2}$ invariante y $e_i'=\gamma e_i$}; al pasar $x$ a centímetros, la pendiente se divide por $100$ (es el cambio de $y$ por centímetro, no por metro). $\blacksquare$
\comprobacion la predicción es equivariante: transformar los datos y ajustar da lo mismo que ajustar y transformar la predicción.
\end{solucion}

\begin{solucion}{eje:t6-regresion-media}
\planteamiento escribir la recta como $\widehat y-\ybar=\widehat b_1(x-\xbar)$ con $\widehat b_1=r\,s_y/s_x$.
\pasos
\begin{enumerate}[label=\arabic*.]
  \item (a) $\widehat y-\ybar=r\dfrac{s_y}{s_x}(x-\xbar)$; dividiendo por $s_y$: $\dfrac{\widehat y-\ybar}{s_y}=r\,\dfrac{x-\xbar}{s_x}$, es decir, $\widehat z_y=r\,z_x$.
  \item (b) $\abs{\widehat z_y}=\abs r\,\abs{z_x}<\abs{z_x}$ si $\abs r<1$ y $z_x\ne0$: la predicción está más cerca de la media, en unidades de desviación típica.
  \item (c) $\widehat z_y=0{,}5\cdot2=1$.
\end{enumerate}
\resultado \textbf{$\widehat z_y=r\,z_x$; con $r=0{,}5$, un modelo a $+2$ desviaciones en la primera evaluación se predice a $+1$ en la segunda.} $\blacksquare$
\comprobacion si $\abs r=1$ no hay regresión a la media ($\widehat z_y=\pm z_x$); si $r=0$, la mejor predicción es la media. Es una razón para desconfiar del modelo que encabeza una única evaluación ruidosa.
\end{solucion}
